\documentclass[12pt,a4paper,openright,twoside]{book}


% input e lingua
\usepackage[utf8]{inputenc}
\usepackage[english,italian]{babel} % italiano come lingua principale

% Template e pacchetti del tamplate
\usepackage{disi-thesis}
\usepackage{code-lstlistings}
\usepackage{notes}
\usepackage{shortcuts}

% Per le note:
% Possibili utilizzi: 
% (1) Nota testuale: \ldots\todo{Finire la frase}
% (2) Quando manca immagine: \missingfigure{Questo è semplicemente un promemoria per inserire un'immagine}
% (3) IMPORTANTE: per generare una lista delle note TODO, basta fare \listoftodos
% (4) Alcune opzioni avanzate: \todo[options]{text}
%     - backgroundcolor, textcolor, linecolor, bordercolor -> passati dentro "options"
\usepackage{todonotes}

% Acronomi: usa 'acro' (NON 'acronym')
\usepackage{acro}

% Per spezzare parole lunghe (es. URL, nomi di file, ecc.)
\usepackage{seqsplit}


% --- Pacchetti consigliati (decommenta se necessario) ---
\usepackage{booktabs}
\usepackage{siunitx}
\usepackage{pgfplots}
\pgfplotsset{compat=1.18}
\usepackage{subcaption}
% \usepackage{graphicx}
\usepackage{siunitx}
\usepackage{array}
\sisetup{detect-all,round-mode=places,round-precision=2}

% Per tabelle con larghezza automatica
\usepackage{tabularx,threeparttable}
\newcolumntype{Y}{>{\raggedright\arraybackslash}X}
\usepackage{ragged2e} % Per un migliore allineamento del testo a capo

% --- Valori inseriti direttamente nel testo/tabelle (nessuna macro numerica) ---

% Per i box
\usepackage[many]{tcolorbox}

% Definiamo un nuovo ambiente per i box "Rischi & Mitigazioni" (versione corretta)
\newtcolorbox{riskbox}[2][]{
  colback=black!5!white,
  colframe=black!75!black,
  fonttitle=\bfseries,
  % --- QUESTA E' LA LINEA MODIFICATA ---
  title=Box~\thetcbcounter: #2, % Costruisce il titolo con il contatore automatico
  #1
}

% Opzioni utili (facoltative):
\acsetup{
  make-links = true,              % link dalla lista alle occorrenze
  first-style = long-short,       % prima volta: "long (SHORT)" (oppure short-long)
  list/name = {Elenco degli acronimi}
}

% macro per produrre l'abbreviazione "e.g." -> exempli gratia - "per esempio"
\newcommand{\eg}{\emph{e.g.}\space}

% Bibliografia: usa biblatex+biber per \parencite
\usepackage{csquotes}
\usepackage[
  backend=biber,
  % style=alphabetic,   % cambia se preferisci
  style=authoryear,
  sorting=ynt
]{biblatex}
\addbibresource{bibliography.bib}

\school{\unibo}
\programme{Corso di Laurea in Ingegneria e Scienze Informatiche}
\title{Sviluppo di un Sistema di Visione Artificiale per l'Analisi di Video di Pallacanestro}
\author{Nome Cognome}
\date{\today}
\subject{Visione Artificiale}
\supervisor{Prof. Nome Cognome}
\session{II}
\academicyear{2024-2025}

% Definition of acronyms
% Acronyms (acro)
\DeclareAcronym{iot}{ short = IoT , long = Internet of Things }
\DeclareAcronym{vm}{ short = VM , long = Virtual Machine , short-plural = VMs , long-plural = Virtual Machines }
\DeclareAcronym{od}{ short = OD , long = Object Detection }
\DeclareAcronym{ot}{ short = OT , long = Object Tracking }
\DeclareAcronym{cv}{ short = CV , long = Computer Vision }
\DeclareAcronym{rgb}{ short = RGB , long = Red Green Blue }
\DeclareAcronym{yuv}{ short = YUV , long = Luminance (Y) and Chrominance (UV) Color Space }
\DeclareAcronym{fps}{ short = FPS , long = frames per second }
\DeclareAcronym{id}{ short = ID , long = identità }
\DeclareAcronym{kpi}{ short = KPI , long = Key Performance Indicator }
\DeclareAcronym{kpis}{ short = KPIs , long = Key Performance Indicators }
\DeclareAcronym{ai}{ short = AI , long = Artificial Intelligence }
\DeclareAcronym{roi}{ short = ROI , long = Region of Interest }
\DeclareAcronym{kf}{ short = KF , long = Kalman Filter }

% --- Nuovi (dal tuo blocco glossaries) ---
\DeclareAcronym{nms}{  short = NMS ,  long = non-maximum suppression }
\DeclareAcronym{iou}{  short = IoU ,  long = intersection over union }
\DeclareAcronym{ap}{  short = AP ,  long = average precision }
\DeclareAcronym{map}{  short = mAP ,  long = mean average precision }
\DeclareAcronym{mota}{ short = MOTA , long = multiple object tracking accuracy }
\DeclareAcronym{mot}{ short = MOT , long = multiple object tracking }
\DeclareAcronym{motp}{ short = MOTP , long = multiple object tracking precision }
\DeclareAcronym{f1}{ short = F1 , long = F1-score }
\DeclareAcronym{idf1}{ short = IDF1 , long = ID F1-score }
\DeclareAcronym{hota}{ short = HOTA , long = higher order tracking accuracy }
\DeclareAcronym{cpu}{  short = CPU ,  long = central processing unit , short-plural-form = CPUs , long-plural-form = central processing units }
\DeclareAcronym{gpu}{  short = GPU ,  long = graphics processing unit , short-plural-form = GPUs , long-plural-form = graphics processing units }
\DeclareAcronym{reid}{ short = Re-ID , long = Re-Identification }
\DeclareAcronym{rpn}{ short = RPN , long = Region Proposal Network }
\DeclareAcronym{fpn}{ short = FPN , long = Feature Pyramid Network }
\DeclareAcronym{detr}{ short = DETR , long = Detection Transformer }
\DeclareAcronym{cnn}{ short = CNN , long = Convolutional Neural Network }
\DeclareAcronym{rnn}{ short = RNN , long = Recurrent Neural Network }
\DeclareAcronym{yolo}{ short = YOLO, long  = You Only Look Once }
\DeclareAcronym{coco}{ short = COCO, long  = Common Objects in Context }
\DeclareAcronym{json}{ short = JSON, long  = JavaScript Object Notation }
\DeclareAcronym{qa}{ short = QA , long  = Quality Assurance }
\DeclareAcronym{mse}{ short = MSE , long  = Mean Squared Error }
\DeclareAcronym{ram}{ short = RAM, long  = random-access memory }
\DeclareAcronym{vram}{ short = VRAM, long  = video random-access memory }
\DeclareAcronym{onnx}{ short = ONNX, long = Open Neural Network Exchange }
\DeclareAcronym{mlops}{ short = MLOps, long = Machine Learning Operations }
\DeclareAcronym{mlflow}{ short = MLflow, long = Machine Learning Flow }
\DeclareAcronym{io}{ short = I/O, long = Input/Output }
\DeclareAcronym{cuda}{ short = CUDA, long = Compute Unified Device Architecture }
\DeclareAcronym{rmse}{ short = RMSE, long = Root Mean Squared Error }
\DeclareAcronym{idsw}{ short = IDSW , long = identity switch }
\DeclareAcronym{ml}{ short = ML , long = Mostly Lost }
\DeclareAcronym{mt}{ short = MT , long = Machine Tracked }
\DeclareAcronym{mae}{ short = MAE , long = Mean Absolute Error }
\DeclareAcronym{px}{ short = px , long = pixel }
\DeclareAcronym{m}{ short = m , long = metri }
\DeclareAcronym{ade}{ short = ADE , long = Average Displacement Error }
\DeclareAcronym{fde}{ short = FDE , long = Final Displacement Error }
\DeclareAcronym{ssd}{ short = SSD , long = Single Shot MultiBox Detector }
\DeclareAcronym{r-cnn}{ short = R-CNN , long = Region-based Convolutional Neural Network }
\DeclareAcronym{sort}{ short = SORT , long = Simple Online and Realtime Tracking }
\DeclareAcronym{oc-sort}{ short = OC-SORT , long = Observation-Centric Simple Online and Realtime Tracking }
\DeclareAcronym{hr-net}{ short = HR-NET , long = High-Resolution Network }
\DeclareAcronym{clip}{ short = CLIP , long = Contrastive Language-Image Pre-training }
\DeclareAcronym{tta}{ short = TTA , long = Test-Time Augmentation }
\DeclareAcronym{csv}{ short = CSV , long = Comma-Separated Values }
\DeclareAcronym{fiba}{ short = FIBA , long = Fédération Internationale de Basketball }
\DeclareAcronym{dlt}{ short = DLT , long = Direct Linear Transformation }
\DeclareAcronym{ransac}{ short = RANSAC , long = Random Sample Consensus }
\DeclareAcronym{ema}{ short = EMA , long = Exponential Moving Average }
\DeclareAcronym{lmeds}{ short = LMEDS , long = Least Median of Squares }








% #########################################################################################
% COME USARLI:

% Prima occorrenza → "long (SHORT)", poi solo "SHORT"
% Nel campo della \ac{cv} rientrano \ac{od} e \ac{ot}.

% Maiuscola a inizio frase
% \Ac{od} è spesso valutata con \ac{iou} e \ac{map}.

% Plurale automatico (usa le chiavi *-plural se servono forme speciali)
% Abbiamo eseguito test su due \acp{vm} con una \ac{gpu}, ottenendo 30 \ac{fps}.

% Forme esplicite quando servono
% La \aclong{cpu} è diversa dalla \acshort{gpu}.
% #########################################################################################


\mainlinespacing{1.241} % line spacing in mainmatter, comment to default (1)

\begin{document}

\frontmatter\frontispiece

\begin{abstract}	
Max 2000 characters, strict.
\end{abstract}

\begin{dedication} % this is optional
\begin{quote}
\emph{``I may not have gone where I intended to go, but I think I have ended up where I needed to be.''} 
\\[0.5em]
\guillemotleft \emph{Forse non sono andato dove intendevo andare, ma credo di essere arrivato dove dovevo essere.}\guillemotright 
\\
--- Douglas Adams, \emph{The Long Dark Tea-Time of the Soul} (1988)
\end{quote}
\end{dedication}

%----------------------------------------------------------------------------------------
\tableofcontents % Indice
\printacronyms % Dove vuoi la lista (tipicamente dopo l'indice o prima del corpo)
\listoffigures     % (optional) comment if empty
\lstlistoflistings % (optional) comment if empty
\listoftodos

%----------------------------------------------------------------------------------------

\mainmatter

% THIS WAS INITIALLY INSIDE THE INTRODUCTION CHAPTER - IT MAY BE USEFUL DO NOT DELETE IT - |FOR NOW|
% You can use acronyms that your defined previously,
% such as \ac{iot}.
% %
% If you use acronyms twice,
% they will be written in full only once
% (indeed, you can mention the \ac{iot} now without it being fully explained).
% %
% In some cases, you may need a plural form of the acronym.
% %
% For instance,
% that you are discussing \ac{vms},
% you may need both \ac{vm} and \ac{vms}.

% \paragraph{Structure of the Thesis}

% \note{At the end, describe the structure of the paper}
%----------------------------------------------------------------------------------------
\chapter{Introduzione (2-3 pagine)}
\label{chap:intro}
%----------------------------------------------------------------------------------------
L'analisi dei dati ha trasformato radicalmente il mondo dello sport,
sia professionistico sia amatoriale. Lungi dall'essere un dominio
affidato esclusivamente all'intuizione di allenatori e atleti, lo
sport moderno è diventato un terreno fertile per l'applicazione di
metodologie quantitative mirate a ottenere un vantaggio competitivo.
In questo scenario, la pallacanestro si distingue per la sua
natura dinamica, la rapidità delle azioni e l'alta densità di eventi
significativi, che la rendono al contempo un soggetto ideale e
complesso per l'analisi automatizzata.

Questo capitolo introduce il lavoro di tesi, partendo dal contesto
e dalle motivazioni che lo hanno ispirato. Verranno quindi definiti
in modo formale i problemi affrontati e gli obiettivi specifici
del progetto. Infine, verrà delineata la struttura del documento
per guidare il lettore attraverso le diverse fasi del lavoro.

% ++++++++++++++++++++++++++++++++++++++++++++++++++++++++++++++++++++++++++++++++++++++
% Richiesta:
% Iniziare descrivendo la crescente importanza dell'analisi dati e della tecnologia nel mondo dello sport, 
% con un focus sul basket. Spiegare come l'analisi automatizzata possa fornire insight preziosi a coach, 
% atleti e analisti, superando i limiti di tempo e costo dell'analisi manuale.

\section{Contesto e Motivazioni}
\label{sec:contesto-motivazioni}
Negli ultimi anni, l'analisi dei dati nello sport (talvolta indicata
come \emph{sports analytics}) ha visto una crescita esponenziale,
alimentata dalla disponibilità di grandi quantità di dati, 
in gran parte sotto forma di registrazioni video. L'analisi
manuale di queste ore di filmati, svolta da video analisti, è un
processo dispendioso in termini di tempo, costoso e
intrinsecamente soggettivo. Da qui l'esigenza di sviluppare
sistemi automatici capaci di estrarre informazioni oggettive e ad
alto valore aggiunto direttamente dai video.

La visione artificiale (in inglese \ac{cv}) è una branca
dell'intelligenza artificiale (\ac{ai}) che mira a replicare la capacità
umana di comprendere e interpretare il mondo visivo, offrendo gli
strumenti per affrontare questa sfida. Applicata al basket, essa
promette di automatizzare il rilevamento (\ac{od}) dei giocatori,
il tracciamento (\ac{ot}) della palla e l'identificazione di eventi di gioco,
fornendo a coach, atleti e scout dati utili per l'analisi
tattica, la valutazione delle prestazioni e lo sviluppo tecnico.

Il contesto della pallacanestro indoor presenta tuttavia sfide
specifiche: la velocità del gioco genera motion blur; le frequenti
occlusioni, con giocatori che si sovrappongono o nascondono la palla,
complicano il tracciamento; le condizioni di illuminazione artificiale
possono creare riflessi sul parquet; le divise delle squadre possono
avere colori simili, rendendo difficile la distinzione.

Questa tesi si inserisce in tale contesto con l'obiettivo di
progettare e sviluppare un sistema prototipale che, partendo da un
video di una partita di basket, sia in grado di eseguire due compiti
fondamentali per qualsiasi analisi successiva: il conteggio
dei giocatori presenti in campo e la ricostruzione della
traiettoria della palla. Il contributo pratico consiste nel fornire
una pipeline software modulare e documentata che affronta queste sfide
utilizzando tecniche allo stato dell'arte, aprendo la strada ad analisi
più complesse.
% ++++++++++++++++++++++++++++++++++++++++++++++++++++++++++++++++++++++++++++++++++++++

% ++++++++++++++++++++++++++++++++++++++++++++++++++++++++++++++++++++++++++++++++++++++
% Richiesta:
% Definire in modo chiaro e formale i problemi che la tesi affronta:
% \begin{enumerate}
%   \item La necessità di un sistema automatico per ricostruire la traiettoria della palla durante un tiro 
%   a canestro a partire da un video.
%   \item L'esigenza di contare in modo affidabile il numero di giocatori presenti in una determinata area 
%   del campo o in un frame video.
%   \item ?Altri?
% \end{enumerate}
\section{Il Problema}
\label{sec:problema}
Sulla base del contesto delineato, il problema generale può essere
scomposto in due sotto-problemi computazionali principali, come
illustrato schematicamente in Figura~\ref{fig:scenario_e_compiti}.

\begin{figure}[h]
  \centering
  \includegraphics[width=\textwidth]{figures/01_intro/PrimaFiguraAlternativa3-CAP1.png}
  % \framebox[1.0\textwidth][c]{Schema a blocchi dell'architettura 
  % concettuale del sistema}
  \caption{Scenario e compiti della tesi. Il sistema riceve in
  input un video di una partita e, attraverso moduli di analisi
  specializzati, produce dati strutturati (conteggio giocatori,
  traiettoria palla) e visualizzazioni. Le principali sfide sono
  rappresentate dalle occlusioni, dall'elevata velocità degli 
  oggetti e dalle condizioni di illuminazione.}
  \label{fig:scenario_e_compiti}
\end{figure}

\begin{enumerate}
  \item \textbf{Rilevamento e conteggio dei giocatori.} Dato un video $V$ composto
  da una sequenza di frame $\{f_1, f_2, \ldots, f_N\}$, il primo
  problema consiste nello stimare, per ogni frame $f_t$, il numero
  totale $c_t$ di giocatori presenti all'interno dell'inquadratura o
  in una specifica regione di interesse (\ac{roi}).
  Il sistema deve essere resistente a occlusioni parziali,
  all'affollamento di giocatori in aree ristrette (ad esempio l'area dei tre secondi) 
  e alla presenza di non-giocatori dall'aspetto simile (ad esempio arbitri o pubblico 
  vicino al campo).
  \item \textbf{Rilevamento e stima della traiettoria della palla.} Il secondo
  problema consiste nello stimare, per ogni frame $f_t$ in cui la
  palla è in movimento, la sua posizione $p_t \in \mathbb{R}^2$ sul piano immagine. 
  L'algoritmo deve gestire sfide intrinseche: le piccole dimensioni della palla, 
  l'alta velocità (con conseguente motion blur) e, soprattutto, le frequenti 
  e prolungate occlusioni da parte dei giocatori. 
  Una stima accurata della traiettoria è prerequisito per il rilevamento di 
  eventi come tiri, passaggi e rimbalzi.
\end{enumerate}

Questi problemi pratici sollevano alcuni quesiti di ricerca che
la tesi si propone di investigare:
\begin{itemize}
  \item Qual è il compromesso ottimale tra accuratezza e velocità di
  inferenza per i modelli di rilevamento di oggetti (\ac{od}) applicati 
  al dominio specifico del basket?
  \item In che misura un approccio ibrido, che combini un rilevatore
  basato su deep learning con un filtro di stato classico (Filtro di
  Kalman -- \ac{kf}), può migliorare la continuità del tracciamento della palla
  in presenza di occlusioni?
  \item Quali strategie di \emph{data augmentation} risultano più efficaci 
  per migliorare la capacità di generalizzazione dei modelli in un
  ambiente controllato ma complesso come un campo da basket indoor?
\end{itemize}
% ++++++++++++++++++++++++++++++++++++++++++++++++++++++++++++++++++++++++++++++++++++++

% ++++++++++++++++++++++++++++++++++++++++++++++++++++++++++++++++++++++++++++++++++++++
% Richiesta:
% Elencare gli scopi concreti del progetto:
% \begin{itemize}
%   \item Studiare lo stato dell'arte delle tecniche di computer vision per l'analisi del basket.
%   \item Creare e/o adattare un dataset di video, etichettandolo per creare un ground truth affidabile.
%   \item Sviluppare e addestrare un modello per il rilevamento dei giocatori.
%   \item Implementare un algoritmo per il tracciamento della traiettoria della palla.
%   \item Valutare le performance dei modelli sviluppati tramite metriche quantitative.
% \end{itemize}
\section{Obiettivi della Tesi}
\label{sec:obiettivi-tesi}

Per rispondere ai quesiti di ricerca e affrontare i problemi
definiti, sono stati stabiliti i seguenti obiettivi concreti, la
cui mappatura con le metriche di valutazione e i capitoli di
riferimento è riassunta in Tabella~\ref{tab:obiettivi_metriche_capitoli}.

\paragraph{Revisione critica dello stato dell'arte.}
Analizzare la letteratura scientifica per identificare 
le tecniche più promettenti per il rilevamento e il tracciamento 
di oggetti in contesti sportivi, con un focus specifico sulla 
pallacanestro.

\begin{figure}[htbp]
  \centering
  \includegraphics[width=\textwidth]{figures/01_intro/SecondaFiguraAlternativa-CAP1.png}
  \caption{Ambito e assunzioni del progetto.}
  \label{fig:ambito-assunzioni}
\end{figure}

\paragraph*{Assunzioni operative}
\begin{itemize}[leftmargin=*, itemsep=0.2em]
  \item \textbf{Scenario:} partite di basket indoor.
  \item \textbf{Telecamera:} una telecamera fissa (necessaria per l'omografia).
  \item \textbf{Qualità video:} risoluzione $\geq$\,720p; frame rate 30 \ac{fps}.
  \item \textbf{Privacy:} nessuna identificazione; i volti non sono processati né salvati.
  \item \textbf{Limiti funzionali:} nessuna analisi tattica complessa e nessun 
  riconoscimento individuale.
\end{itemize}

\begin{table}[h!]
  \centering
  \caption{Mappatura degli obiettivi della tesi con le relative metriche di valutazione 
  e i capitoli di riferimento.}
  \label{tab:obiettivi_metriche_capitoli}
  \begin{tabularx}{\textwidth}{l Y l}
    \toprule
    \textbf{Obiettivo} & \textbf{Metrica di Valutazione / Output} & \textbf{Capitolo} \\
    \midrule
    Revisione dello stato dell'arte & Sintesi critica della letteratura 
    & Cap.~\ref{chap:stato-dell-arte} \\
    \addlinespace
    Creazione e annotazione del dataset & Dataset con \emph{ground truth} 
    (bounding box di giocatori e palla) & Cap.~\ref{chap:metodologia-sviluppo} \\
    \addlinespace
    Sviluppo del modello di rilevamento & Precision, Recall, \ac{map} 
    & Cap.~\ref{chap:risultati-sperimentali} \\
    \addlinespace
    Tracciamento giocatori & \ac{idf1}; \ac{hota}; 
    (di supporto: \ac{mota}/\ac{motp}, \acf{idsw}, \acf{mt}/\acf{ml}) 
    & Cap.~\ref{chap:risultati-sperimentali} \\
    \addlinespace
    Traiettoria palla & \acf{mae} su piano omografico in \acf{px} o \acf{m}; 
    \ac{rmse}; \acf{ade}/\acf{fde} & Cap.~\ref{chap:risultati-sperimentali} \\
    \addlinespace
    Valutazione complessiva & Analisi quantitativa e qualitativa dei risultati 
    ottenuti & Cap.~\ref{chap:conclusioni-sviluppi-futuri} \\
    \bottomrule
  \end{tabularx}
\end{table}
% ++++++++++++++++++++++++++++++++++++++++++++++++++++++++++++++++++++++++++++++++++++++

% ++++++++++++++++++++++++++++++++++++++++++++++++++++++++++++++++++++++++++++++++++++++
% Richiesta: 
% Fornire una breve panoramica dei capitoli successivi, guidando il lettore attraverso il documento.
\section{Struttura della Tesi}
\label{sec:struttura-tesi}
Il presente documento è organizzato come segue.

\paragraph{Capitolo~\ref{chap:stato-dell-arte} -- Stato dell'Arte.} 
Introduce i concetti fondamentali della visione artificiale e 
analizza i lavori scientifici e i dataset pubblici più 
rilevanti per l'analisi del basket.

\paragraph{Capitolo~\ref{chap:metodologia-sviluppo} -- Metodologia e Sviluppo del Progetto.} 
Descrive in dettaglio l'architettura del sistema, 
la preparazione dei dati e le scelte implementative 
per ciascun modulo di analisi.

\paragraph{Capitolo~\ref{chap:risultati-sperimentali} -- Risultati Sperimentali e Valutazione.} 
Presenta i risultati quantitativi ottenuti, includendo studi 
di ablazione, e discute criticamente le prestazioni e i limiti 
del sistema.

\paragraph{Capitolo~\ref{chap:conclusioni-sviluppi-futuri} -- Conclusioni e Sviluppi Futuri.}
Riassume il lavoro svolto, ne ribadisce i contributi e delinea 
possibili direzioni per estensioni e miglioramenti futuri.
% ++++++++++++++++++++++++++++++++++++++++++++++++++++++++++++++++++++++++++++++++++++++

%----------------------------------------------------------------------------------------
\chapter{Stato dell'Arte (10--15~pagine)}
\label{chap:stato-dell-arte}
%----------------------------------------------------------------------------------------
\noindent Questo capitolo sintetizza i principali concetti, metodi 
e risultati della letteratura rilevanti per la progettazione di un 
sistema di visione artificiale applicato a video di pallacanestro indoor.
L'obiettivo è fornire un quadro coerente e aggiornato che consenta di 
motivare le scelte metodologiche del Capitolo~\ref{chap:metodologia-sviluppo} 
e di definire criteri di valutazione oggettivi per il 
Capitolo~\ref{chap:risultati-sperimentali}.

% ++++++++++++++++++++++++++++++++++++++++++++++++++++++++++++++++++++++++++++++++++++++
\section{Introduzione alla Visione Artificiale}
\label{sec:introduzione-vision-artificiale}

La visione artificiale (\ac{cv}) studia algoritmi e sistemi capaci 
di estrarre informazioni semantiche da dati visivi, tipicamente immagini e video. 
I compiti più pertinenti a questa tesi comprendono: rilevamento di oggetti
(\ac{od}: localizzare e classificare istanze con \emph{bounding box});
tracciamento di oggetti (\ac{ot}: stimare \ac{id} e traiettorie nel tempo);
stima di posa (\emph{pose estimation}: coordinate articolari 2D/3D del corpo umano); 
re-identificazione (\ac{reid}: associazione della stessa identità attraverso frame o videocamere); 
riconoscimento di eventi (\emph{event recognition}: classificare o segmentare azioni 
complesse nel tempo, ad esempio tiro, rimbalzo, ``pick-and-roll'').

Queste componenti possono essere integrate in una pipeline che, 
a partire dal video, produce indicatori utili ad allenatori e analisti, 
come conteggi dei giocatori, traiettorie della palla, mappe di tiro o 
rilevamento di situazioni tattiche ricorrenti.

Nel contesto della \textbf{pallacanestro indoor}, la letteratura 
evidenzia vincoli peculiari: occlusioni frequenti dovute alla 
densità di giocatori in area; illuminazione non uniforme e 
riflessi del parquet che alterano colori e contrasto; 
\emph{\emph{camera shake}} e \emph{panning} nelle riprese 
\emph{handheld} o da tribuna; la palla è piccola e si muove 
rapidamente, con \emph{motion blur}; uniformi simili e la 
presenza di arbitri/pubblico possono confondere i modelli. 
Tali fattori influenzano sia l'accuratezza 
(\emph{precision}/\emph{recall}, \emph{\ac{map}}, \emph{\ac{idf1}}) 
sia l'efficienza (\emph{\ac{fps}}, \emph{ms/frame}, memoria), 
e guidano scelte di \emph{compute budget} 
(\emph{\ac{cpu}}/\emph{\ac{gpu}}) e di \emph{data strategy} 
(dataset, annotazioni, \emph{data augmentation}). 
Nei paragrafi seguenti si presentano gli elementi 
fondamentali (Sezione~\ref{sec:elementi-visione-artificiale}), 
le famiglie di \emph{detector} e \emph{tracker} più diffuse 
(Sezione~\ref{sec:algoritmi-rilevamento-tracciamento}), 
applicazioni al basket (Sezione~\ref{sec:cv-applicata-basket})
e dataset pubblici utili (Sezione~\ref{sec:dataset-pubblici}).
% sezioni e label verificate e allineate
% ++++++++++++++++++++++++++++++++++++++++++++++++++++++++++++++++++++++++++++++++++++++


% ++++++++++++++++++++++++++++++++++++++++++++++++++++++++++++++++++++++++++++++++++++++
% Richiesta:
% Fornire una breve introduzione ai concetti fondamentali necessari per comprendere il progetto: elaborazione
% di immagini, object detection, object tracking.
\section{Elementi di Visione Artificiale}
\label{sec:elementi-visione-artificiale}

Questa sezione introduce, con definizioni sintetiche, i costrutti 
necessari a comprendere le scelte progettuali adottate 
successivamente.

\paragraph{Rappresentazioni e pre-elaborazione.} Un frame video è una 
matrice di intensità in uno spazio colore (\emph{\ac{rgb}}; 
\emph{\ac{yuv}}). Operazioni comuni includono ridimensionamento, 
normalizzazione, bilanciamento del bianco, riduzione del rumore 
e \emph{deblurring}. La frequenza dei fotogrammi (\ac{fps}) 
e la risoluzione (ad esempio $1920\times1080$) influenzano 
latenza e accuratezza.

\paragraph{Rilevamento di oggetti (\ac{od}).} Dato un frame, l'\ac{od} 
produce un insieme di \emph{bounding box} con etichette (ad esempio 
\texttt{giocatore} e \texttt{palla}). Le metriche standard sono \emph{\ac{iou}},
sovrapposizione \emph{box--ground truth} e \ac{map}, media delle 
\ac{ap} su soglie di \ac{iou}. La \emph{\ac{nms}} rimuove box ridondanti.

\paragraph{Tracciamento di oggetti (\ac{ot}).} 
L'\ac{ot} associa le detezioni nel tempo per stimare 
\ac{id} persistenti e traiettorie. Si distingue tra 
\emph{tracking-by-detection} (associazione 
tra detezioni \emph{frame-to-frame}) e 
\emph{joint detection/tracking} 
(rilevamento e associazione unificati). 
Metriche: \emph{\ac{mota}} e \emph{\ac{motp}} 
(CLEAR \emph{\ac{mot}} \cite{Bernardin2008CLEAR}), 
\ac{idf1} \cite{Ristani2016Performance}, \emph{\ac{hota}} 
\cite{Luiten2020HOTA}.

\paragraph{Stima di posa e \ac{reid}.} La stima di posa (\emph{pose}) 
fornisce punti chiave (\emph{keypoints}) scheletrici; utile per 
distinguere posture (tiro, sottomano, rimbalzo).
Esempi: \emph{OpenPose} \cite{8765346}, \emph{HRNet} \cite{Sun2019Deep}.
La re-identificazione (\ac{reid}) apprende descrittori di aspetto 
per riassociare la stessa persona; è spesso addestrata su dataset 
come \emph{Market-1501} \cite{7410490}.

\paragraph{Efficienza e \emph{compute budget}.} La latenza
(ms/frame), il \emph{throughput} (\ac{fps}) e la memoria
determinano la praticabilità in tempo reale. Il tempo
dipende da: modello (profondità, risoluzione di input, 
\emph{backbone}), hardware (solo \ac{cpu} vs \ac{gpu}; 
ad esempio \ac{gpu} \emph{consumer} o \emph{embedded}), 
\emph{batch size}. In ambienti indoor con \emph{camera shake} 
e \emph{motion blur}, \emph{temporal smoothing} 
(filtri di Kalman, medie esponenziali) e strategie di 
\emph{debouncing} possono stabilizzare \ac{id} e traiettorie.

\paragraph{Vincoli specifici del basket indoor.} 
Occlusioni: frequenti in area; richiedono stabilità nell' 
\emph{\ac{idsw}}. Illuminazione/riflessi: il parquet lucido introduce 
aberrazioni; la palla arancione può confondersi con toni caldi. 
Palla piccola/veloce: serve risoluzione sufficiente e/o 
\emph{temporal cues}. Uniformi simili e arbitri/pubblico: 
necessarie strategie di \emph{class imbalance handling} e 
\emph{data augmentation} mirata.

\subsection*{Punti chiave}
\begin{itemize}
  \item \ac{od}, \ac{ot}, \emph{pose} e \ac{reid} sono i mattoni 
  principali; metriche chiave: \ac{iou}, \ac{map}, \ac{idf1}, \ac{hota}.
  \item Latenza e memoria dipendono da risoluzione, \ac{fps} e hardware; 
  la qualità video influenza direttamente le prestazioni.
  \item I vincoli del basket (occlusioni, riflessi, palla veloce) 
  richiedono strategie specifiche di acquisizione e \emph{pre-processing}.
\end{itemize}

\subsection*{Implicazioni per il Cap.~\ref{chap:metodologia-sviluppo}.}
Giustificare scelte di risoluzione (ad esempio $1280\times720$ vs 
$1920\times1080$), un target di 30~\ac{fps}, e del \emph{compute budget} 
(solo \ac{cpu} vs \ac{gpu}). Prevedere un modulo di stabilizzazione temporale 
e \emph{data augmentation} specifica per parquet/uniformi.
% ++++++++++++++++++++++++++++++++++++++++++++++++++++++++++++++++++++++++++++++++++++++

% ++++++++++++++++++++++++++++++++++++++++++++++++++++++++++++++++++++++++++++++++++++++
\section{Algoritmi di rilevamento e tracciamento di oggetti}
\label{sec:algoritmi-rilevamento-tracciamento}
Presentiamo le famiglie di metodi più adottate, con particolare 
enfasi sulla stabilità in ambienti indoor e sul rapporto costo/efficacia.

\subsection{Rilevamento: due-stadi, uno-stadio e transformer}

\paragraph{Due-stadi (two-stage).} I detector \emph{two-stage} 
generano \emph{region proposals} e raffinano classificazione/localizzazione.
\emph{Faster \ac{r-cnn}} riduce il collo di bottiglia tramite una 
\emph{Region Proposal Network} \cite{Ren2015FasterRCNN}.
\textbf{Vantaggi}: alta accuratezza su oggetti medi/grandi e a \ac{iou} elevate; 
\textbf{Svantaggi}: latenza superiore e sensibilità alla risoluzione.

\paragraph{Uno-stadio (one-stage).} I detector \emph{one-stage} 
apprendono predizioni dense: \emph{\ac{ssd}} \cite{Liu2016SSD} e le famiglie
\emph{\ac{yolo}} \cite{Redmon2018YOLOv3}. \textbf{Vantaggi}: 
\emph{throughput} elevato, buone prestazioni su \emph{hardware commodity}; 
\textbf{Svantaggi}: potenziale calo su oggetti piccoli ad \emph{iou} alte. 
Varianti recenti (ad esempio architetture \emph{anchor-free} 
\cite{Ultralytics2025AnchorFree} e \emph{decoupled head} 
\cite{Meng2023YOLOv5sFog}) hanno migliorato il compromesso 
\emph{accuracy/latency}. Nel presente progetto si adotta un detector 
\emph{one-stage} (\emph{Ultralytics \ac{yolo}} 
\cite{Ge2021YOLOX}) per giocatori e palla. 
% \todo{aggiungere citazioni puntuali per anchor-free 
% (ad esempio FCOS/ATSS) e decoupled head (ad esempio YOLOX).}

\paragraph{Basati su transformer.} \emph{\ac{detr}} riformula 
il rilevamento come \emph{set prediction} con meccanismi di 
attenzione \cite{Ye2022OSTrack}. Versioni successive 
(\emph{Deformable \ac{detr}} \cite{Zhu2020DeformableDETR}) accelerano
la convergenza e migliorano il rilevamento di oggetti piccoli. 
\textbf{Vantaggi}: qualità stabile e associazione naturale tra box e oggetti; 
\textbf{Svantaggi}: addestramento più complesso e possibile costo computazionale maggiore.

\paragraph{Pratiche per oggetti piccoli/veloci.} Per la palla: 
input ad alta risoluzione; \emph{feature pyramid} (\ac{fpn})~\cite{Lin2017Feature, Liu2018Path};
\emph{focal loss} per classi rare; \emph{temporal priors} 
(propagazione delle detezioni). Per i giocatori: \emph{soft-\ac{nms}} e
\emph{\ac{iou}-aware heads} per mitigare \emph{crowding} e sovrapposizioni. 
Nel progetto si adottano inferenze localizzate su \emph{\ac{roi}} e su 
tasselli (\emph{tiles}) ad alta risoluzione per incrementare il 
\emph{recall} della palla, con selezione dei candidati tramite fusione 
di punteggi visivi e coerenza dinamica (\emph{gating} basato su filtro di Kalman).

\subsection{Tracciamento multi-oggetto (\ac{mot})}

\paragraph{Tracking-by-detection.} \emph{\ac{sort}} utilizza un filtro di 
Kalman e l'assegnazione dei costi con l'algoritmo Hungarian \cite{Kuhn1955HungarianMethod, 
Munkres1957Assignment}, offrendo velocità molto alta ma sensibilità alle occlusioni. 
% \todo{aggiungere citazione per SORT (Bewley et al., 2016) e, se desiderato, 
% per Hungarian (Kuhn 1955/Munkres 1957).}
\emph{DeepSORT} aggiunge \ac{reid} per stabilizzare le identità in presenza 
di \emph{occlusions} \cite{Wojke2017DeepSort}. 
\emph{ByteTrack} enfatizza l'uso di detezioni a bassa confidenza per mantenere le 
tracce durante cali temporanei \cite{Wang2021TransformerTracker}.
\emph{\ac{oc-sort}} sfrutta l'\emph{order consistency} per gestire in modo efficace il 
\emph{camera motion} \cite{Ye2022OSTrack}. 
\emph{CenterTrack} integra rilevamento e spostamento centrato sul keypoint 
dell'oggetto \cite{Zhou2020TrackingObjects}.

\paragraph{Joint detection/tracking e transformer.} Metodi recenti unificano le 
due fasi: \emph{TrackFormer} sfrutta query persistenti per le tracce 
\cite{Meinhardt2021TrackFormer}. Questi approcci possono ridurre i 
\emph{cambi di \ac{id}} (\emph{ID switches}) in \emph{crowded scenes}, 
al prezzo di addestramenti più onerosi e maggiore domanda di dati annotati.

\paragraph{Metriche per \ac{mot}.} \ac{mota} (errori globali: 
falsi positivi/negativi e cambi di \ac{id})~\cite{Braso2020Learning}, \ac{motp} (precisione 
media di localizzazione), \ac{idf1} (F1 sulle identità corrette) e 
\ac{hota} (bilancia accuratezza di localizzazione e associazione) sono 
standard \cite{Bernardin2008CLEAR, Ristani2016Performance, Luiten2020HOTA}. 
In fase di progettazione è utile riportare anche la latenza 
(ms/frame), gli \ac{fps} su hardware specifico (\ac{cpu}/\ac{gpu}, ad esempio 
\ac{gpu} mobile vs \emph{desktop}), la memoria e la risoluzione. 
Ogni numero va accompagnato da \emph{dataset} e \emph{split} 
(ad esempio \texttt{MOT17 test}) per confronti riproducibili.

\subsection{Pose e \ac{reid} nel contesto sportivo}

\paragraph{Pose.} \emph{OpenPose} (multi-persona, parti del corpo come 
mappe di confidenza) \cite{Cao2016Realtime} e \emph{\ac{hr-net}} (rappresentazioni 
ad alta risoluzione) \cite{Sun2019Deep} sono riferimenti consolidati. 
Nella pallacanestro, la posa fornisce indizi per rimbalzi, \emph{close-out} 
e preparazione al tiro; tuttavia occlusioni e \emph{motion blur} possono ridurre 
la qualità dei \emph{keypoint}.

\paragraph{\ac{reid}.} \emph{DeepSORT} è un esempio di integrazione leggera. 
Una \ac{reid} specifica per team può sfruttare le informazioni cromatiche 
della divisa, \emph{texture} delle maglie e numeri, ma richiede annotazioni 
curate; modelli generici addestrati su \emph{Market-1501} \cite{Zheng2015Scalable} 
possono necessitare di \emph{fine-tuning} per ambienti indoor.

\subsection*{Punti chiave}
\begin{itemize}
  \item \emph{Two-stage} (accuratezza) vs \emph{one-stage} (latenza): 
  scegliere in funzione del \emph{compute budget} e della dimensione 
  degli oggetti.
  \item Nel \ac{mot}, l'aggiunta di \ac{reid} (DeepSORT) o l'associazione 
  a bassa soglia (ByteTrack) migliora la stabilità in scenari affollati.
  \item Pose e \ac{reid} aggiungono segnali utili ma introducono costo 
  computazionale e richiedono dati annotati.
\end{itemize}

\paragraph{Implicazioni per il Cap.~\ref{chap:metodologia-sviluppo}.}
Se l'obiettivo è il \emph{real-time} su \ac{gpu} \emph{consumer}, adottare 
un \emph{detector one-stage} (Ultralytics \ac{yolo}) e il tracciamento integrato 
con identità persistenti per i giocatori; per la palla, combinare un \emph{detector} 
specializzato con ricerca localizzata (\ac{roi}/\emph{tile}) e filtro di Kalman per 
la coerenza temporale. Prevedere opzionalmente la \emph{Pose} per analisi di eventi, 
con \emph{fallback} disattivabile su hardware limitato. 


\todo{da scaletta - valutare se lasciare o rimuovere}
\subsection{Approcci Tradizionali (cenni):}
\label{sec:approcci-tradizionali}
Menzionare brevemente tecniche classiche (ad esempio \emph{background subtraction}) 
e i loro limiti in contesti dinamici come lo sport.

\todo{da scaletta - valutare se lasciare o rimuovere}
\subsection{Approcci basati su Deep Learning:}
\label{sec:approcci-deep-learning}
Concentrarsi sulle Reti Neurali Convoluzionali (\ac{cnn}). Descrivere architetture 
moderne per l'\emph{object detection} come \ac{yolo}, \ac{ssd} o Faster \ac{r-cnn}, 
spiegando perché sono particolarmente adatte al contesto sportivo. 
Per il tracking, introdurre concetti come i filtri di Kalman o \emph{tracker} più moderni 
come DeepSORT.
% ++++++++++++++++++++++++++++++++++++++++++++++++++++++++++++++++++++++++++++++++++++++


% ++++++++++++++++++++++++++++++++++++++++++++++++++++++++++++++++++++++++++++++++++++++
% Richiesta:
% Analizzare la letteratura scientifica esistente. Esaminare cosa è già stato fatto in questo campo, cercando
% studi su:
% \begin{itemize}
%   \item Tracciamento dei giocatori (player tracking).
%   \item Rilevamento di eventi (tiro, palleggio, passaggio).
%   \item Tracciamento della palla.
%   \item Analisi tattica automatizzata.
% \end{itemize}
\section{Visione artificiale applicata alla pallacanestro}
\label{sec:cv-applicata-basket}

La letteratura sulla pallacanestro combina rilevamento/tracciamento 
di giocatori e palla, stima di eventi e analisi tattica.
Rispetto ad altri sport, il basket presenta \emph{crowding} elevato 
e spazi ridotti. Si distinguono lavori su (i) \emph{player tracking}, 
(ii) \emph{ball tracking}, (iii) \emph{event detection}, 
(iv) analisi tattica multi-agente.

\subsection{Player tracking}
Approcci \emph{tracking-by-detection} con \emph{detector} rapidi e associazione 
robusta sono prevalenti: combinazioni di \emph{detector} tipo 
\ac{yolo} e/o \ac{ssd} con DeepSORT/ByteTrack sono comuni in scenari indoor grazie 
al buon compromesso accuratezza/latenza. In presenza di uniformi simili e 
\emph{camera motion}, l'aggiunta di \ac{reid} e vincoli spaziali (modelli di 
moto con filtro di Kalman) riduce i \emph{cambi di \ac{id}}. L'uso della
\emph{pose} aiuta a stabilizzare il \emph{matching} durante \emph{partial occlusions}.
% \todo{aggiungere citazioni specifiche su player tracking nel basket, 
% ad esempio studi su dataset indoor e sistemi \ac{mot} per basket.}
% \todo{Esempio placeholder (mantieni le citazioni per 
% compatibilità se non usi biblatex): 
% \cite{todo-player-tracking-basket}}

\paragraph{Collegamento al progetto.} Nel progetto, il tracciamento dei giocatori 
è ottenuto con Ultralytics \ac{yolo} tramite la funzione di \emph{tracking} con identità 
persistenti; le identità fornite sono mantenute \emph{frame-to-frame} e utilizzate 
per l'assegnazione della squadra con un classificatore visuale basato su \ac{clip} 
\cite{Chen2024ContrastiveLocalized} che opera sui ritagli dei giocatori e mantiene 
coerenza temporale.

\subsection{Ball detection/tracking}
La palla è piccola (pochi pixel a risoluzioni standard), veloce e soggetta 
a occlusioni. Metodi efficaci combinano: (i) detector specializzati per 
oggetti piccoli (\ac{fpn}, input ad alta risoluzione) \cite{Chen2009Physics}, 
(ii) \emph{temporal priors} con filtri di Kalman o particellari, 
(iii) \emph{motion cues} (flusso ottico) per distinguere \emph{motion blur} 
da falsi positivi. Alcuni lavori sfruttano \emph{multi-view} con omografie 
del campo per triangolazione quando sono disponibili più camere.
% \todo{aggiungere citazioni specifiche su ball tracking nel basket: \cite{todo-ball-tracking-basket}}

\paragraph{Collegamento al progetto.} Il tracciamento della palla combina 
un detector dedicato con: (i) ricerca localizzata su \ac{roi} attorno 
all'ultima posizione nota e su tasselli ad alta risoluzione, (ii) fusione 
di punteggi visivi (confidenza del detector, indizi di colore e forma) e 
spaziali (prossimità ai giocatori, coerenza con posizione/velocità previste),
(iii) filtro di Kalman per continuità e recupero durante le occlusioni; 
soglie e parametri si adattano dinamicamente quando la palla è persa. 
E' disponibile, opzionalmente, la \emph{\ac{tta}} \cite{Kimura2024Understanding}.


\subsection{Rilevamento di eventi}
Gli eventi (tiro, rimbalzo, assist, palla persa) richiedono modellazione 
temporale oltre il singolo frame. Architetture \emph{\ac{cnn}+\ac{rnn}} 
o \emph{transformer} temporali sono impiegate per \emph{action recognition} 
\cite{Xiao2023Basketball} in ambito sportivo; nel basket si combinano 
traiettorie di palla/giocatori, pose e, quando disponibile, segnali audio. 
La disponibilità di annotazioni temporali fini limita l'addestramento 
supervisionato; strategie \emph{weakly supervised} \cite{Zhu2023Weaker} 
e \emph{self-supervised} \cite{Gui2023Survey} sono state esplorate.
% \todo{aggiungere citazioni su event detection nel basket (supervisioni 
% deboli/auto-supervisionate): \cite{todo-event-detection-basket}}

\paragraph{Collegamento al progetto.} Il rilevamento del tiro è 
realizzato con regole deterministiche: \emph{trigger} quando la palla
sale rapidamente e non è in possesso; conclusione quando scende sotto 
l'altezza del canestro; l'esito deriva dalla distanza minima della 
traiettoria rispetto all'area del canestro. La classificazione 
2/3 punti usa la distanza del tiratore dal canestro in vista tattica o, 
in assenza, un \emph{fallback} 2D~\cite{White2020Analysis}.


\subsection{Analisi tattica}
Una volta acquisite traiettorie e identità, si derivano 
\emph{\ac{kpis}} come spaziature, velocità medie, \emph{heatmap} 
di presenza e sequenze ricorrenti (\emph{set plays})~\cite{Mikolajec2013Game}. 
L'inferenza di ruoli e schemi è facilitata dalla \emph{pose} 
e dalle posizioni relative con vincoli di geometria del campo
(omografia del campo)~\cite{Lingrui2025TacticExpert}.
% \todo{aggiungere citazioni su analisi tattica automatica nel basket: 
% \cite{todo-analisi-tattica-basket}}

\paragraph{Collegamento al progetto.} Le posizioni dei giocatori 
sono mappate su una vista tattica 2D mediante omografia del campo;
si proiettano i canestri nel frame, si validano i \emph{keypoint} del 
campo e si ottengono coordinate normalizzate per metriche come
spaziature, velocità e distanze, oltre a \emph{overlay} dedicati.


\subsection{Vincoli pratici indoor e mitigazioni}
\begin{itemize}
  \item \textbf{Occlusioni prolungate}: usare \emph{tracklet stitching} 
  e \ac{reid}; sfruttare \emph{camera elevation} quando possibile.
  \item \textbf{Illuminazione/riflessi}: \emph{color constancy} e 
  \emph{data augmentation} (variazioni di luminanza/saturazione); 
  \emph{polarizing filter} in acquisizione quando consentito.
  \item \textbf{Palla piccola/veloce}: input a maggiore risoluzione 
  su \emph{crop} dinamici attorno a zone salienti; \emph{temporal votes}.
  \item \textbf{Uniformi simili/arbitri/pubblico}: classificatori di 
  ruolo/squadra basati su \emph{color histograms} stabilizzati e numeri 
  di maglia quando leggibili.
  \item \textbf{Camera shake}: stima del \emph{global motion} e 
  stabilizzazione video; \emph{debouncing} sulle \acp{id}.
\end{itemize}


\subsection*{Punti chiave}
\begin{itemize}
  \item Il \emph{player tracking} e il \emph{ball tracking} 
  richiedono strategie diverse: dimensione e dinamica della 
  palla impongono moduli dedicati.
  \item Annotazioni temporali e di identità rappresentano 
  un collo di bottiglia per \emph{event detection} e analisi tattica.
  \item I vincoli indoor impongono scelte di acquisizione e 
  \emph{pre-processing} tanto quanto di modello.
\end{itemize}

\subsection*{Implicazioni per il Cap.~\ref{chap:metodologia-sviluppo}}
Prevedere una pipeline modulare: detector per i giocatori 
con tracking integrato (Ultralytics \ac{yolo} con identità persistenti)
e un ramo specializzato per la palla che combini detector, ricerca 
localizzata (\ac{roi}/tile) e filtro di Kalman. Integrare, dove possibile, 
l'omografia del campo per normalizzare le coordinate.
% ++++++++++++++++++++++++++++++++++++++++++++++++++++++++++++++++++++++++++++++++++++++


% ++++++++++++++++++++++++++++++++++++++++++++++++++++++++++++++++++++++++++++++++++++++
% Richiesta:
% Descrivere brevemente i dataset già e spiegare perché, nonostante la loro esistenza, si è reso necessario 
% registrare e/o etichettare video specifici per questo progetto.
\section{Dataset pubblici per l'analisi sportiva}
\label{sec:dataset-pubblici}

I dataset influenzano in modo decisivo la trasferibilità dei risultati. 
E' utile distinguere dataset generici per \ac{od} e \ac{mot} da dataset 
specifici per lo sport. Quando si riportano le prestazioni è essenziale 
indicare \emph{split} (\textbf{train/val/test}), risoluzione, \ac{fps}, 
hardware e metriche.

\subsection{Dataset generici per \ac{od} e \ac{mot}}

\paragraph{\ac{coco}.} MS \ac{coco} è un riferimento per l'\ac{od} 
con annotazioni dense su molte classi \cite{Lin2014Microsoft}. 
Il pre-addestramento su \ac{coco} migliora la generalizzazione 
alla classe \texttt{person} (giocatori).
% TODO: forse avrebbe senso dire che offre già di base, per esempio, 
% anche la palla sportiva (non specificamente quella da basket però)
% e altre cose che potrebbero essere utili da introdurre per la 
% comprensione del lettore. 

\paragraph{MOTChallenge.} \emph{MOT16/17/20} sono \emph{benchmark} 
per il \ac{mot} in scene affollate; forniscono annotazioni 
\texttt{person} e metriche standard (\ac{mota}, \ac{idf1}, \ac{hota}). 
Questi dataset abilitano il confronto tra tracker e la valutazione 
delle prestazioni in condizioni di 
\emph{crowding}~\cite{Milan2016MOT16, Leal-Taixe2017MOT17, Dendorfer2020MOT20}.
% \todo{aggiungere citazioni per MOT16/17/20.}

\paragraph{Pose e \ac{reid}.} Per la \emph{pose}: 
\ac{coco} Keypoints (sottoinsieme di \ac{coco})~\cite{Lin2014Microsoft}.
Per la \emph{re-identificazione} (\ac{reid}): \emph{Market-1501}~\cite{Zheng2015Scalable}.

\subsection{Dataset nel dominio sportivo con focus sulla pallacanestro}

\paragraph{Dataset sportivi multi-disciplina.} \emph{SoccerNet}~\cite{Giancola2018Soccernet} 
è ampio e include annotazioni temporali per eventi calcistici; 
pur non essendo cestistico, fornisce buone pratiche per annotazioni 
di eventi e valutazioni con allineamento temporale. Altri benchmark
multi-sportivi come \emph{ActivityNet}~\cite{CabaHeilbron2015ActivityNet}
sono utili per apprendere feature generali sul movimento umano.

\paragraph{Dataset dedicati alla pallacanestro.} Alcuni dataset includono 
\emph{player/ball tracking}, spesso multi-camera con 
omografie del campo, ma non sempre sono pubblici o 
completi di annotazioni a livello di 
evento~\cite{dOrey2020APIDIS, Karun2022DeepSport, Pan2023BasketballInstants}.

\subsection{Linee guida di valutazione}

\begin{itemize}
  \item \textbf{Detection (\texttt{giocatori}, \texttt{palla})}: 
  riportare \ac{map} a soglie di \ac{iou} (ad esempio 0.5 e 0.5:0.95), 
  specificando risoluzione di input e latenza (ms/frame) su hardware 
  (\ac{gpu}/\ac{cpu}).
  \item \textbf{Tracking (\ac{mot})}: \ac{map}, \ac{motp}, 
  \ac{idf1}, \ac{hota} su \emph{split} dichiarati (ad esempio 
  \texttt{val} interno, \texttt{test} pubblico), indicando 
  \ac{fps} e memoria.
  \item \textbf{Traiettoria della palla}: errore medio per frame su 
  coordinate normalizzate al campo; specificare \ac{fps},
  frame scartati ammessi e strategia di interpolazione.
  \item \textbf{Eventi}: F1 per classe e media pesata; 
  finestra di tolleranza per l'allineamento temporale esplicitata.
\end{itemize}

\paragraph{Collegamento al progetto.} La valutazione include 
curve \emph{precision-recall} e una stima di \ac{ap} per il
rilevamento di palla e giocatori, oltre a metriche di continuità 
della palla (\emph{coverage}, lunghezza media/massima delle
sequenze di \emph{miss}/\emph{hit}, ritardo della prima corretta 
detezione); i risultati vengono anche esportati come grafici e 
file \emph{\ac{csv}}.

\subsection{Motivazioni per un dataset ad hoc}
Le risorse pubbliche esistenti risultano spesso insufficienti 
per il contesto di pallacanestro indoor per:
\begin{itemize}
  \item \textbf{Mismatch di etichette}: molti dataset annotano 
  solo \texttt{person}, senza palla o ruoli di squadra.
  \item \textbf{Assenza di eventi}: mancano annotazioni temporali 
  fini (ad esempio inizio/fine di tiro, rimbalzo) utili all'analisi.
  \item \textbf{Dominio visivo diverso}: illuminazione indoor, 
  riflessi del parquet, uniformi simili e \emph{camera motion} 
  differiscono dai dataset generici.
  \item \textbf{Vincoli multi-camera e omografia}: i meta-dati 
  geometrici non sono sempre disponibili.
  \item \textbf{Licenze/uso}: alcune collezioni non consentono 
  rilascio o condivisione di clip/annotazioni derivate.
\end{itemize}
Questi fattori motivano la registrazione e l'etichettatura 
di un dataset mirato, con protocolli di split riproducibili e meta-dati completi.


\subsection*{Punti chiave}
\begin{itemize}
  \item \ac{coco} e MOTChallenge restano riferimenti per 
  pre-addestramento e confronto base sulla classe 
  \texttt{person}.
  \item La disponibilità di dataset cestistici 
  completi è limitata; spesso è necessario creare 
  o adattare un dataset proprietario.
  \item Report oggettivi richiedono la specifica di split, 
  risoluzione, \ac{fps}, hardware e metriche.
\end{itemize}

\subsection*{Implicazioni per il Cap.~\ref{chap:metodologia-sviluppo}}
Pianificare la costruzione di un dataset \emph{ad hoc} 
con annotazioni per giocatori e palla, protocolli di 
split riproducibili e meta-dati (risoluzione, \ac{fps}, 
setup camera). Prevedere \emph{transfer learning} da 
\ac{coco} e validazione con metriche \ac{mot}.
% ++++++++++++++++++++++++++++++++++++++++++++++++++++++++++++++++++++++++++++++++++++++

%----------------------------------------------------------------------------------------
\chapter{Metodologia e Sviluppo del Progetto (15-25 pagine)}
\label{chap:metodologia-sviluppo}
%----------------------------------------------------------------------------------------
Questo capitolo descrive in dettaglio la metodologia adottata 
e le scelte implementative che hanno guidato lo sviluppo 
del sistema di analisi video per la pallacanestro. 
L'obiettivo è fornire una panoramica chiara e riproducibile 
del lavoro svolto, collegando le decisioni tecniche alle 
necessità pratiche dell'analisi sportiva e alle evidenze emerse dallo 
stato dell'arte -- Cap.~\ref{chap:stato-dell-arte}.

Partendo dall'architettura generale del sistema, verranno 
analizzati i singoli componenti, dalla preparazione dei dati 
ai modelli di rilevamento (object detection) e tracciamento 
(object tracking), fino alle tecniche utilizzate per estrarre 
metriche quantitative e visualizzazioni significative.
Ogni scelta sarà giustificata nel contesto dei vincoli 
operativi, come le condizioni di illuminazione indoor, 
le occlusioni e la velocità degli oggetti di interesse.

% ++++++++++++++++++++++++++++++++++++++++++++++++++++++++++++++++++++++++++++++++++++++
% Richiesta:
% Descrivere il flusso di lavoro complessivo del progetto, possibilmente con un diagramma. 
% Esempio: Input Video $\rightarrow$ Pre-elaborazione Frame $\rightarrow$ Modello di Rilevamento 
% Giocatori $\rightarrow$ Modello di Tracciamento Palla $\rightarrow$ Output Dati Analitici.
\section{Architettura del Sistema}
\label{sec:architettura-sistema}

L'efficacia del sistema risiede in un'architettura modulare e 
sequenziale, progettata per trasformare un video grezzo di una 
partita di pallacanestro in un insieme strutturato di dati analitici 
e annotazioni grafiche. Ogni modulo della pipeline è specializzato in un 
compito specifico e propaga il proprio output al modulo successivo, 
creando un flusso di lavoro logico e facilmente ispezionabile.

\subsection{Scenario d'uso e Flusso End-to-End}

Lo scenario d'uso primario prevede l'analisi di video 
registrati da una singola telecamera fissa, posizionata 
in modo da inquadrare l'intero campo o una sua porzione 
significativa. La stabilità della telecamera è un prerequisito 
fondamentale per il corretto funzionamento di alcuni moduli, 
in particolare per la stima dell'omografia del campo.

L'input del sistema è un file video (ad esempio in formato 
\texttt{.mp4} o \texttt{.MOV}).
L'output è un nuovo file video annotato in cui sono sovrimpresse 
informazioni quali:
\begin{itemize}
  \item i riquadri di delimitazione (bounding box) e 
  gli \ac{id} univoci dei giocatori;
  \item l'assegnazione di ciascun giocatore a una squadra;
  \item la traiettoria della palla;
  \item una vista tattica 2D che mappa la posizione dei 
  giocatori su un diagramma del campo;
  \item metriche calcolate in tempo reale, come la velocità 
  e la distanza percorsa dai giocatori.
\end{itemize}

Il flusso di elaborazione end-to-end (Figura~\ref{fig:architettura}) 
è orchestrato come segue:

\begin{enumerate}
  \item \textbf{Acquisizione e pre-elaborazione del video}: il 
  sistema acquisisce il video sorgente e lo decodifica in una 
  sequenza di frame individuali. Ogni frame rappresenta 
  un'istantanea temporale della partita che verrà analizzata 
  singolarmente.
  \item \textbf{Rilevamento dei punti chiave e calcolo dell'omografia del campo}: 
  per ogni frame, un modello di deep learning rileva i punti 
  chiave del campo (ad esempio linea del tiro libero, centro 
  del campo, linee di fondo, ecc.). Questi punti sono usati 
  per calcolare una matrice di omografia, ovvero una 
  trasformazione proiettiva che mappa le coordinate di ogni 
  punto dal piano dell'immagine del video a un piano 2D 
  standardizzato del campo da basket.
  \item \textbf{Rilevamento e tracciamento dei giocatori}: 
  un modello di \ac{od}, basato sull'architettura 
  \ac{yolo}v8 e addestrato specificamente su immagini 
  di giocatori di basket, identifica la posizione di 
  ogni giocatore nel frame. Gli output del rilevatore 
  (bounding box) vengono passati a un algoritmo di 
  tracciamento che assegna un \ac{id} univoco a ogni 
  giocatore e ne mantiene l'identità attraverso i 
  frame successivi.
  \item \textbf{Assegnazione delle squadre}: un modulo analizza 
  il colore dominante all'interno del bounding box di ogni 
  giocatore tracciato per assegnarlo a una delle due squadre. 
  Questo permette di distinguere tra compagni e avversari 
  nelle analisi successive.
  \item \textbf{Rilevamento e tracciamento della palla}: un 
  rilevatore \ac{yolo}v8 specializzato e un filtro di Kalman 
  lavorano in sinergia per localizzare e tracciare la palla. 
  Il filtro di Kalman è cruciale per gestire le occlusioni 
  (quando la palla è nascosta da un giocatore) e la 
  sfocatura da movimento (motion blur), predicendo la 
  posizione della palla sulla base del suo stato precedente.
  \item \textbf{Conversione in vista tattica}: utilizzando 
  la matrice di omografia calcolata al punto 2, le coordinate 
  dei giocatori e della palla vengono trasformate e proiettate 
  su un'immagine statica del campo da basket, generando una 
  vista tattica 2D.
  \item \textbf{Calcolo delle metriche}: moduli dedicati 
  utilizzano i dati di tracciamento per calcolare metriche 
  quantitative, come la velocità istantanea e la distanza 
  totale percorsa da ogni giocatore.
  \item \textbf{Visualizzazione e generazione dell'output}: 
  i risultati di tutti i moduli precedenti vengono utilizzati 
  per arricchire visivamente il frame originale. 
  Vengono disegnati i bounding box, le traiettorie, le 
  informazioni sulla velocità e la vista tattica. 
  I frame annotati vengono infine ricodificati in un nuovo file video.
\end{enumerate}

% \missingfigure{Manca architettura end-to-end del sistema. Il diagramma mostra il flusso sequenziale dei dati, 
% dall'acquisizione del video sorgente fino alla generazione dell'output annotato. Ogni blocco rappresenta un modulo 
% funzionale specializzato.}
\begin{figure}[h]
  \centering
  \includegraphics[width=\textwidth]{figures/01_intro/FlowchartpipelineE2E-B.png}
  % \framebox[1.0\textwidth][c]{Diagramma a blocchi dell'architettura end-to-end}
  \caption{Architettura end-to-end del sistema. Il diagramma 
  mostra il flusso sequenziale dei dati, dall'acquisizione 
  del video sorgente fino alla generazione dell'output annotato. 
  Ogni blocco rappresenta un modulo funzionale specializzato.}
  \label{fig:architettura}
\end{figure}

\subsection{Componenti Tecnologici Chiave}

La realizzazione di questa architettura si basa su un insieme 
di librerie e framework open-source consolidati nell'ambito 
della \ac{cv}:
\begin{itemize}
  \item \textbf{Python 3.10}: linguaggio di programmazione 
  principale per l'intero progetto.
  \item \textbf{OpenCV}: utilizzato per tutte le operazioni 
  di base di elaborazione delle immagini e per la gestione 
  dei flussi video (lettura e scrittura).
  \item \textbf{Ultralytics--YOLOv8}: framework di deep learning 
  impiegato per l'addestramento e l'inferenza dei modelli di 
  rilevamento per giocatori, palla e punti chiave del campo.
  \item \textbf{NumPy}: libreria fondamentale per la manipolazione 
  efficiente di array e matrici, utilizzata in quasi tutti i moduli.
  \item \textbf{FilterPy}: libreria che fornisce un'implementazione 
  del filtro di Kalman, utilizzata nel modulo di tracciamento della 
  palla.
  \todo{forse dovrei rimuoverlo perché non lo uso più?}
\end{itemize}

\subsection*{Punti chiave}
\begin{itemize}
  \item L'architettura è modulare e sequenziale, facilitando lo 
  sviluppo, il testing e il debug dei singoli componenti.
  \item Il sistema combina approcci basati su deep learning 
  (\ac{yolo}v8 per il rilevamento) con algoritmi classici di 
  visione artificiale (filtro di Kalman, calcolo dell'omografia).
  \item La stima dell'omografia è un passaggio chiave che abilita 
  la generazione della vista tattica, un potente strumento di analisi.
\end{itemize}

\subsection*{Implicazioni per il Cap.~\ref{chap:risultati-sperimentali}}
\begin{itemize}
  \item La valutazione delle performance dovrà essere condotta sia 
  a livello di singolo modulo (ad esempio accuratezza del rilevatore 
  di giocatori con metrica \ac{map}) sia a livello end-to-end 
  (ad esempio latenza complessiva misurata in frame al secondo, \ac{fps}).
  \item L'efficacia della vista tattica dipenderà direttamente 
  dalla precisione del modello di rilevamento dei punti chiave del 
  campo, la cui performance dovrà essere quantificata.
  \item L'accuratezza del tracciamento (ad esempio con metriche 
  come \ac{mota} o \ac{idf1}) sarà un indicatore fondamentale della 
  capacità del sistema di mantenere l'identità dei giocatori nel tempo.
\end{itemize}
% ++++++++++++++++++++++++++++++++++++++++++++++++++++++++++++++++++++++++++++++++++++++


% ++++++++++++++++++++++++++++++++++++++++++++++++++++++++++++++++++++++++++++++++++++++

\section{Raccolta e Preparazione dei Dati}
\label{sec:raccolta-dati}

La qualità e la pertinenza dei dati sono il fondamento su cui poggia l'efficacia di qualsiasi sistema di visione artificiale. Per questo progetto è stata adottata una \textbf{strategia di acquisizione dati multi-componente}, volta a costruire un corpus di addestramento robusto e a validarlo in scenari realistici. 

L'approccio si è articolato in tre fasi principali:
\begin{enumerate}
  \item \textbf{Utilizzo di dataset pubblici specializzati} 
  come solide basi di partenza per i diversi compiti di rilevamento.
  \item \textbf{Estensione e arricchimento di tali dataset} 
  attraverso un meticoloso processo di etichettatura manuale, 
  che ha permesso di aumentarne la varietà e la complessità.
  \item \textbf{Utilizzo di video personalizzati e online} 
  sia come fonte di nuovi dati di training, sia come set di 
  test indipendente per valutare la generalizzazione del sistema.
\end{enumerate}

Questa metodologia ha permesso di beneficiare di ampie basi di 
dati già annotate e, al contempo, di creare risorse su misura 
per le finalità del progetto, garantendo una valutazione finale 
approfondita e affidabile.

\subsection{Dataset Utilizzati}
\label{sec:dataset-utilizzati}

\noindent Per specializzare e potenziare un modello pre-addestrato 
(\ac{yolo}v8) sul dominio specifico del basket, è stato necessario un 
processo di addestramento mirato (\emph{fine-tuning}). A tale scopo, 
sono stati realizzati due dataset personalizzati, partendo da risorse 
pubbliche e arricchendole tramite un processo strutturato di etichettatura 
manuale. Le fonti e la composizione di tali dataset sono descritte di seguito.

\begin{enumerate}
  \item \textbf{Fonti video per l'arricchimento dei dataset}:
  Il cuore del lavoro di personalizzazione dei dati è consistito 
  nell'analisi e nell'annotazione di fotogrammi estratti da 
  un'ampia raccolta di video. Per garantire una grande varietà di 
  scenari, sono state utilizzate due tipologie di fonti video, 
  le cui immagini sono state impiegate per arricchire 
  \textbf{entrambi} i dataset descritti ai punti successivi:
  \begin{itemize}
    \item \textbf{Video personalizzati}: Una raccolta di 
    filmati registrati dall'autore con una \emph{telecamera 
    consumer} (smartphone) a 1080p e 30~fps. Lo scopo era 
    catturare scenari amatoriali (ad esempio campetti 
    all'aperto e allenamenti) e prospettive non comuni nelle 
    riprese professionali.
    \item \textbf{Video da fonti online}: Una selezione di 
    registrazioni di partite di basket (ad esempio, 
    highlights da YouTube) per includere dati con 
    inquadrature televisive e un'alta intensità di gioco, 
    tipici di contesti professionali.
  \end{itemize}
  Per entrambe le fonti, è stato impiegato uno \emph{script} Python 
  per automatizzare il campionamento dei fotogrammi 
  (frame sampling), estraendone sistematicamente uno a un 
  intervallo regolare di 0.4 secondi. L'intero set di 
  fotogrammi così ottenuto è stato poi sottoposto al processo 
  di etichettatura manuale. I video personalizzati, nella loro 
  interezza, sono stati inoltre usati per l'analisi qualitativa 
  finale del sistema.


  \item \textbf{Dataset per il rilevamento di giocatori e azioni di gioco}:
  Come base di partenza è stato impiegato il dataset pubblico 
  \texttt{Basketball Players Computer Vision Model}~\cite{Basketball-Player-Detection-2_Dataset}. 
  Questa risorsa, già annotata con bounding box per 
  \textbf{11 classi} (entità e azioni di gioco), è stata 
  ampliata integrando i fotogrammi etichettati manualmente 
  provenienti dalle fonti video descritte al punto precedente. 
  La versione finale del dataset così 
  ottenuto~\cite{Basketball-Player-Detection-2-68vvt_Dataset} è 
  stata poi aumentata artificialmente (\texttt{aug3x}) e 
  configurata disabilitando l'augmentation ``mosaic'' 
  (\texttt{nomosaic}), una scelta metodologica volta a evitare 
  la creazione di scene poco realistiche nel contesto sportivo.

  \item \textbf{Dataset per il rilevamento dei punti chiave del campo}:
  Il punto di partenza per questo compito è stato il dataset 
  pubblico \texttt{Basketball Court Detection}~\cite{Basketball-Court-Detection-2_Dataset}, 
  contenente immagini annotate con punti chiave (keypoints) dei 
  riferimenti geometrici del campo. Analogamente al caso 
  precedente, questo dataset è stato arricchito con i fotogrammi 
  etichettati manualmente e provenienti dalla medesima raccolta 
  di video personali e online. Il risultato di questo lavoro di 
  integrazione è il dataset finale impiegato nel 
  progetto~\cite{Basketball-Court-Detection-2-Ny1gt_Dataset}.
\end{enumerate}

La composizione e la suddivisione quantitativa dei dataset finali 
impiegati sono riassunte in Tabella~\ref{tab:dataset_split}.

\begin{table}[htbp]
  \centering
  \caption{Riepilogo quantitativo dei dataset finali e della 
  loro suddivisione (split).}
  \label{tab:dataset_split}
  \begin{tabularx}{\textwidth}{ Y r r r r }
    \toprule
    \textbf{Dataset Utilizzato} & \textbf{Totale Immagini} & 
    \textbf{Train Set} & \textbf{Valid Set} & \textbf{Test Set} \\
    \midrule
    
    Giocatori e Azioni \cite{Basketball-Player-Detection-2-68vvt_Dataset} 
    & 4286
    & 3831 (89\%)
    & 253 (6\%)
    & 202 (5\%) \\
    
    \addlinespace
    
    Punti Chiave Campo \cite{Basketball-Court-Detection-2-Ny1gt_Dataset}
    & 2525
    & 2211 (88\%)
    & 180 (7\%)
    & 134 (5\%) \\

    \bottomrule
  \end{tabularx}
\end{table}

% Spiegare come sono stati etichettati i video per ottenere i dati di riferimento. Descrivere il software
% utilizzato, il formato delle etichette e il processo seguito per garantire la qualità delle annotazioni.
\subsection{Annotazione e Creazione del Ground Truth}
\label{sec:annotazione-ground-truth}

\noindent Il processo di arricchimento dei dataset di base, 
descritto nella sezione precedente, si è fondato su un 
meticoloso lavoro di etichettatura manuale. Sebbene i dataset 
di partenza fornissero già un primo set di annotazioni, è 
stato necessario etichettare \emph{ex novo} l'intero 
insieme di fotogrammi campionati dalle fonti video 
(personali e online) per creare un ground truth coerente 
e di alta qualità per i due compiti specifici del progetto.

\begin{itemize}
  \item \textbf{Schemi di Annotazione}: Invece di singole classi, 
  sono stati seguiti due schemi di annotazione distinti, uno per 
  ciascun dataset finale:
  \begin{enumerate}
    \item \textbf{Rilevamento Giocatori e Azioni}: Per questo compito 
    è stato adottato lo schema multi-classe del dataset di partenza, 
    composto da \textbf{11 classi} differenti. Un elenco completo e 
    la loro distribuzione sono visibili in Figura~\ref{fig:players-actions-classes}. 
    Queste includono sia entità fisiche (ad esempio \texttt{player}, 
    \texttt{ball}, \texttt{rim}) sia azioni di gioco specifiche 
    (ad esempio \texttt{player-dribble}, \texttt{player-layup}).
      
    \item \textbf{Rilevamento Punti Chiave}: Per la stima della 
    geometria del campo, è stato definito un template strutturato 
    (o ``scheletro'') composto da \textbf{41 punti chiave}. 
    Come illustrato in Figura~\ref{fig:keypoints-template}, 
    questi punti sono stati posizionati su riferimenti geometrici 
    precisi (angoli delle aree, intersezioni delle linee, estremità 
    del campo) e collegati da segmenti per mappare l'intera struttura 
    del campo da gioco. E' importante sottolineare come, in ogni frame, 
    venissero etichettati unicamente i punti dello scheletro 
    effettivamente visibili nell'inquadratura, gestendo così in 
    modo nativo le viste parziali del campo.
  \end{enumerate}
\end{itemize}

\begin{itemize}
  \item \textbf{Strumento e processo di annotazione}: 
  L'etichettatura è stata eseguita interamente sulla 
  piattaforma \textbf{Roboflow}. La sua interfaccia 
  \emph{web-based} ha permesso di processare i fotogrammi, 
  precedentemente estratti tramite lo script Python, 
  applicando due diverse tipologie di etichette: 
  \textbf{bounding box} per il rilevamento di giocatori e 
  azioni, e \textbf{punti chiave} (keypoints) per la 
  geometria del campo. Un esempio dell'interfaccia durante 
  l'etichettatura con bounding box è mostrato in 
  Figura~\ref{fig:roboflow-ui}.

  \item \textbf{Formato delle etichette}: Per garantire la 
  massima compatibilità, tutte le annotazioni sono state 
  gestite ed esportate in formato \ac{json} secondo lo 
  standard \ac{coco}. Il formato \ac{coco} include le 
  informazioni necessarie all'addestramento (ad esempio 
  coordinate dei bounding box e \ac{id} di classe) ed è 
  ormai uno \emph{standard de facto} per i task di object 
  detection, ampiamente supportato dai principali framework 
  di deep learning, incluso \ac{yolo}v8. Ciò consente di 
  utilizzare i dataset senza ricorrere a conversioni manuali, 
  poiché i framework forniscono utilità per leggere e 
  interpretare le annotazioni.
\end{itemize}

\begin{figure}[h!] % h! per forzare la posizione il più possibile
  \centering
  \includegraphics[width=\textwidth]{figures/03_architettura/roboflow-ui.png} 
  \caption{Esempio dell'interfaccia di annotazione di Roboflow 
  durante l'etichettatura di un frame con bounding box. Sono 
  evidenziati giocatori (magenta), palla (viola), arbitri (ciano), 
  canestro (arancione), numeri delle canottiere (azzurro) e 
  il giocatore che sta tirando (giallo).}
  \label{fig:roboflow-ui}
\end{figure}

\begin{figure}[h!]
  \centering
  \includegraphics[width=\textwidth]{figures/03_architettura/players_actions-classes.png}
  \caption{Elenco delle 11 classi di bounding box definite per il 
  dataset di rilevamento di giocatori e azioni di gioco, con la 
  loro distribuzione numerica e rispettivi colori assegnati.}
  \label{fig:players-actions-classes}
\end{figure}

\begin{figure}[h!]
  \centering
  \includegraphics[width=\textwidth]{figures/03_architettura/keypoints-template2.png}
  \caption{Il template a 41 punti chiave (``scheletro'') utilizzato 
  in Roboflow per l'annotazione della geometria del campo da basket. 
  Sono visibili la numerazione e le connessioni tra i punti.}
  \label{fig:keypoints-template}
\end{figure}

\subsubsection*{Protocollo di \ac{qa}}
\noindent\textbf{Obiettivo.} Garantire coerenza e accuratezza 
delle etichette in un flusso di annotazione svolto da un unico 
operatore.

\smallskip
\noindent\textbf{Strumenti.} Tutte le annotazioni sono state 
effettuate all'interno della piattaforma Roboflow, sfruttandone 
le funzioni integrate di ispezione, analisi statistica e 
versionamento del dataset.

\smallskip
\noindent\textbf{Criteri e controlli adottati.}
\begin{enumerate}[label=(\alph*)]
  \item \textbf{Schema delle classi e nomenclatura.} Definizione a 
  priori degli schemi di annotazione (le 11 classi per i bounding 
  box e lo scheletro a 41 punti) e uso di etichette univoche e 
  consistenti.
  \item \textbf{Policy per i bounding box.} Aderenza stretta ai 
  contorni; esclusione di istanze al di sotto di una soglia minima 
  di area; annotazione della sola porzione visibile in caso di 
  occlusione; gestione delle sovrapposizioni senza duplicazioni.

  \item \textbf{Policy per i punti chiave.} Posizionamento del 
  punto sempre al centro dell'intersezione o nell'angolo più 
  interno della linea di riferimento; coerenza nell'applicazione 
  dello scheletro a 41 punti, anche in caso di prospettive estreme.

  \item \textbf{Linee guida per casi ambigui.} Gestione predefinita 
  di casi difficili come la palla sfocata dal movimento (motion 
  blur) o giocatori quasi totalmente occlusi.

  \item \textbf{Controllo di integrità dei dati (\emph{Sanity check}).} 
  Rimozione di immagini duplicate o corrotte; verifica delle dimensioni 
  e dei rapporti d'aspetto; controllo del bilanciamento tra le classi 
  tramite le statistiche fornite da Roboflow.

  \item \textbf{Suddivisione del dataset.} Suddivisione 
  \emph{train/val/test} stratificata e bloccata per scena, 
  al fine di prevenire contaminazione dei dati (\emph{data leakage}) 
  e preservare la distribuzione statistica in ogni sottoinsieme.

  \item \textbf{Revisione periodica.} Riesame a campione di un 
  sottoinsieme di immagini già annotate per correggere eventuali 
  incoerenze e garantire la stabilità delle scelte nel tempo.

  \item \textbf{Versionamento.} Ogni modifica sostanziale al 
  dataset ha generato una nuova versione tracciata su Roboflow, 
  garantendo la piena riproducibilità degli esperimenti.
\end{enumerate}

\smallskip
\noindent\textbf{Limitazioni.} L'assenza di più annotatori non ha 
permesso di calcolare metriche oggettive come l'accordo inter-annotatore 
(\emph{Inter-Annotator Agreement}). Le revisioni periodiche a campione e i 
controlli automatici della piattaforma sono stati utilizzati come 
principali misure mitigative.

\subsection{Pre-elaborazione e Augmentation}
\label{subsec:augmentation}

\paragraph{Preprocessing (Roboflow).}
Per entrambe le versioni del dataset, prima dell'addestramento sono
stati applicati: \texttt{Auto-Orient} e ridimensionamento con
\emph{stretch} a $640{\times}640$ px. La suddivisione \emph{train/val/test}
è riportata in Tabella~\ref{tab:dataset_split}.

\paragraph{Augmentation offline (generate su Roboflow).}
Per aumentare la variabilità senza introdurre artefatti non presenti
nei dati reali, è stata adottata una triplicazione del set di
addestramento (\texttt{aug3x}: 3 output per immagine). Parametri:
\begin{itemize}
  \item \emph{Giocatori/Azioni}~\cite{Basketball-Player-Detection-2-68vvt_Dataset}:
  \textbf{flip orizzontale}; \textbf{luminosità} in $[-25\%, +25\%]$;
  \textbf{sfocatura} fino a $2$\,px; \textbf{rumore} fino allo $0{,}5\%$
  dei pixel.
  \item \emph{Punti chiave del campo}~\cite{Basketball-Court-Detection-2-Ny1gt_Dataset}:
  (nessun flip); \textbf{luminosità} in $[-15\%, +15\%]$;
  \textbf{sfocatura} fino a $2{,}5$\,px; \textbf{rumore} fino allo
  $0{,}1\%$ dei pixel.
\end{itemize}
\emph{Nota.} Nella generazione \emph{offline} non è stato usato
\emph{mosaic}.

\paragraph{Augmentation online (durante l'addestramento in YOLOv8).}
Durante l'addestramento su Google Colab, oltre ai dati pre-generati su
Roboflow, la procedura di addestramento applica trasformazioni
\emph{online} a ogni mini-batch. Sono stati addestrati due modelli con
configurazioni distinte (parametri essenziali e di augmentation):

\begin{itemize}
  \item \emph{Rilevamento palla (ball-only)} -- \texttt{yolov8m.pt}:\\
  $imgsz=1280$, $epochs=150$, $patience=30$, $batch=-1$ (dimensione
  automatica), $accumulate=4$, $optimizer=\text{auto}$, $lr0=0.003$,
  $lrf=0.2$, $momentum=0.937$, $weight\_decay=5\cdot10^{-4}$,
  $warmup\_epochs=3.0$, $warmup\_momentum=0.8$, $warmup\_bias\_lr=0.1$,
  $rect=\text{True}$, $multi\_scale=\text{False}$;\\
  \emph{augmentation}: $mosaic=0.6$ con chiusura $close\_mosaic=20$,
  $mixup=0.0$, $copy\_paste=0.0$, $hsv\_h=0.015$, $hsv\_s=0.5$,
  $hsv\_v=0.3$, $translate=0.05$, $scale=0.4$, $fliplr=0.2$,
  $flipud=0.0$, $degrees=0.0$, $shear=0.0$, $perspective=0.0$;
  $max\_det=300$.\\
  \textit{Pesi della funzione obiettivo}: $box=8.0$, $cls=0.5$,
  $dfl=1.5$.

  \item \emph{Rilevamento generale (tutte le classi)} -- \texttt{yolov8s.pt}:\\
  $imgsz=960$, $epochs=120$, $patience=30$, $batch=-1$,
  $optimizer=\text{auto}$, $lr0=0.003$, $lrf=0.2$, $momentum=0.937$,
  $weight\_decay=5\cdot10^{-4}$, $warmup\_epochs=3.0$,
  $warmup\_momentum=0.8$, $warmup\_bias\_lr=0.1$, $rect=\text{True}$,
  $multi\_scale=\text{True}$;\\
  \emph{augmentation}: $mosaic=1.0$ con chiusura $close\_mosaic=10$,
  $mixup=0.1$, $copy\_paste=0.0$, $hsv\_h=0.015$, $hsv\_s=0.6$,
  $hsv\_v=0.3$, $translate=0.05$, $scale=0.3$, $fliplr=0.5$,
  $flipud=0.0$, $degrees=0.0$, $shear=0.0$, $perspective=0.0$;
  $max\_det=300$.
\end{itemize}

\paragraph{Scelte e motivazioni (sintesi).}
\begin{itemize}
  \item \textbf{Risoluzione} ($imgsz=1280$ palla; $960$ generale): la
  palla è un target minuto; la risoluzione maggiore migliora la
  localizzazione fine, mentre $960$ bilancia costo e varietà per il
  modello generale.
  \item \textbf{Multi-scala e mosaic}: $multi\_scale$ disattivato per
  la palla per stabilità sul piccolo oggetto; attivato nel generale per
  robustezza a distanza/zoom. \emph{Mosaic} usato solo nelle fasi
  iniziali ($close\_mosaic$) per diversificare senza fissare artefatti.
  \item \textbf{Trasformazioni} (\texttt{hsv}, \texttt{translate},
  \texttt{scale}, \texttt{fliplr}): ampiezze moderate per riflettere
  condizioni indoor realistiche; flip più probabile nel generale per
  incrementare simmetrie, più contenuto sulla palla anche per la
  presenza del flip \emph{offline}.
  \item \textbf{Geometria conservativa} ($degrees$, $shear$,
  $perspective=0$): si evitano distorsioni che altererebbero relazioni
  geometriche del campo e riferimenti utili ai due compiti.
  \item \textbf{Ottimizzazione} (pesi della funzione obiettivo,
  $accumulate$): enfasi sulla componente di localizzazione per la palla
  (\texttt{box}↑, \texttt{cls}↓); accumulo dei gradienti per stabilità
  a risoluzioni elevate.
\end{itemize}


\paragraph{Note di riproducibilità.}
\begin{itemize}
  \item Le augmentation \emph{offline} sono deterministiche per ciascuna
  versione Roboflow citata (\texttt{aug3x}).
  \item Le augmentation \emph{online} sono campionate in fase di
  addestramento secondo i valori riportati sopra; per la replica usare
  identiche impostazioni, $imgsz$ e strategia di chiusura del
  \emph{mosaic}.
  \item Le immagini Roboflow ($640{\times}640$ \emph{stretch}) vengono
  ridimensionate internamente da YOLOv8 a $imgsz\in\{960,1280\}$.
  % Per una riproducibilità completa, specificare (se disponibili)
  % versione esatta di \texttt{ultralytics}, \texttt{torch}, \texttt{cuda}
  % e il \texttt{seed}.
\end{itemize}

Per chiarezza, la Figura~\ref{fig:augmentation} illustra esempi 
delle trasformazioni adottate; l'ultimo pannello mostra il \emph{mosaic}, 
utilizzato esclusivamente in fase di addestramento (augmentation online 
in \ac{yolo}v8).

\begin{figure}[h]
  \centering
  \includegraphics[width=\textwidth]{figures/03_architettura/augmentation_grid.png}
  \caption{Esempi rappresentativi delle tecniche di augmentation adottate.
  Riga superiore: (a) immagine originale; (b) variazione di luminosità (+20\%);
  (c) sfocatura gaussiana (2\,px). Riga inferiore: (d) flip orizzontale (applicato
  solo al dataset giocatori/azioni); (e) rumore (0{,}5\% dei pixel);
  (f) composizione \emph{mosaic}, impiegata esclusivamente in fase di addestramento
  (augmentation online in YOLOv8).}
  \label{fig:augmentation}
\end{figure}

\textbf{Implicazioni per il Cap.~\ref{chap:risultati-sperimentali}.}
\begin{itemize}
  \item La valutazione finale dovrà confrontare le performance del 
  modello sul test set del dataset pubblico e sui video personalizzati, 
  per misurare la capacità di generalizzazione.
  \item Sarà possibile condurre un'analisi degli errori (error analysis) 
  per capire in quali condizioni (ad esempio occlusioni, illuminazione 
  scarsa) il modello fallisce, correlando i fallimenti alle caratteristiche 
  dei dati di training.
  \item Un potenziale studio di ablazione nel Capitolo~\ref{chap:risultati-sperimentali} 
  potrebbe quantificare l'impatto delle diverse tecniche di augmentation 
  sulle performance finali del modello.
\end{itemize}
% ++++++++++++++++++++++++++++++++++++++++++++++++++++++++++++++++++++++++++++++++++++++


% ++++++++++++++++++++++++++++++++++++++++++++++++++++++++++++++++++++++++++++++++++++++
\section{Task 1: Rilevamento e Conteggio dei Giocatori}
\label{sec:task-1}

Il primo compito fondamentale della pipeline di analisi è il 
rilevamento accurato e affidabile di tutti i giocatori presenti 
sul campo. Questo passaggio è il prerequisito per qualsiasi 
analisi successiva, dal tracciamento individuale al 
riconoscimento delle formazioni tattiche. L'obiettivo di 
questo modulo è duplice: identificare la posizione di ogni 
giocatore in ogni frame del video e fornire un conteggio 
numerico degli atleti visibili.

\subsection{Scelta del Modello}
\label{sec:scelta-modello}
% Giustificare la scelta del modello utilizzato (es. "E' stato scelto il modello YOLOv8 per il suo 
% ottimo compromesso tra velocità di inferenza e accuratezza...").

Dopo un'attenta analisi dello stato dell'arte
(Cap.~\ref{chap:stato-dell-arte}) e dei requisiti del progetto,
la scelta è ricaduta sull'architettura \ac{yolo}v8, in quanto
supportata da un insieme di fattori chiave che la rendono
particolarmente adatta al contesto applicativo considerato.

% Dopo un'attenta analisi dello stato dell'arte 
% (Cap.~\ref{chap:stato-dell-arte}) e dei requisiti del 
% progetto, la scelta è ricaduta sull'architettura \ac{yolo} -- 
% versione 8. Questa decisione è motivata da un insieme di fattori 
% chiave che rendono \ac{yolo}v8 particolarmente adatto a questo 
% contesto applicativo.

% CONFRONTA CON QUELLO SOTTO: SONO DUE VERSIONI DELLA STESSA COSA.
% \begin{itemize}
%   \item \textbf{Compromesso accuratezza/velocità}: i modelli della 
%   famiglia YOLO garantiscono un buon equilibrio tra \ac{map} e \ac{fps}. 
%   In analisi video, la velocità di inferenza è un vincolo primario; un 
%   rilevatore a \emph{singolo stadio} come YOLOv8 è quindi preferibile.
%   \item \textbf{Performance consolidate}: al momento dello sviluppo, 
%   YOLOv8 risulta competitivo sui benchmark general-purpose (es. \ac{coco}), 
%   rendendolo una base affidabile per il dominio basket.
%   \item \textbf{Affidabilità su scale diverse}: la variante \emph{anchor-free} 
%   migliora la gestione di oggetti con rapporti d'aspetto e pose variabili 
%   (giocatori in piedi, in salto, piegati, a terra).
%   \item \textbf{Ecosistema}: la catena di strumenti di Ultralytics semplifica 
%   addestramento, validazione ed esecuzione, riducendo il tempo di iterazione.
% \end{itemize}


\begin{itemize}
  \item \textbf{Compromesso tra accuratezza e velocità}: i modelli della 
  famiglia \ac{yolo} offrono un eccellente bilanciamento tra accuratezza 
  di rilevamento (misurata in \ac{map}) e velocità di inferenza (misurata 
  in \ac{fps}). In un'applicazione di analisi video, dove è necessario 
  elaborare decine di frame al secondo, un'alta velocità di inferenza è 
  un requisito non negoziabile. \ac{yolo}v8, essendo un rilevatore a 
  singolo stadio (\emph{single-stage}), eccelle in questo.
  \item \textbf{Performance allo stato dell'arte}: al momento dello 
  sviluppo, \ac{yolo}v8 rappresenta una delle architetture più performanti 
  per il compito di rilevamento (object detection) general-purpose, 
  superando molte alternative in termini di performance su benchmark 
  standard come \ac{coco}.
  \item \textbf{Affidabilità e flessibilità}: l'architettura di \ac{yolo}v8, 
  in particolare nelle varianti più recenti \emph{anchor-free}, mostra 
  una notevole affidabilità nel rilevare oggetti di dimensioni e rapporti 
  d'aspetto variabili, caratteristica utile per identificare giocatori in 
  diverse posture (in piedi, in salto, piegati, a terra, ecc.).
  \item \textbf{Ecosistema e facilità d'uso}: il framework di Ultralytics 
  offre una catena di strumenti completa e intuitiva per 
  l'addestramento, la validazione e la messa in produzione dei modelli, 
  semplificando il ciclo di sviluppo.
\end{itemize}

Le principali sfide del rilevamento nel basket, come l'affollamento 
(alta densità di giocatori in aree ristrette, ad esempio sotto il canestro)
e le occlusioni frequenti e prolungate, sono state considerate. 
Si è ipotizzato che la combinazione di un modello potente come \ac{yolo}v8 
e un dataset aumentato ad hoc (Sezione~\ref{subsec:augmentation}) possa 
fornire la robustezza necessaria per gestire tali scenari.

\begin{table}[h]
  \centering
  \caption{Sintesi delle scelte per il detector dei giocatori e delle 
  principali alternative.}
  \label{tab:scelte_detector}
  \begin{tabularx}{\textwidth}{l X}
    \toprule
    \textbf{Opzione} & \textbf{Motivazione} \\
    \midrule
    \textbf{YOLOv8s} (scelto) & Buon equilibrio accuratezza/velocità; 
    adatto a quasi tempo reale su GPU consumer; ecosistema maturo per 
    addestramento/validazione. \\
    Faster R-CNN (two-stage) & Accuratezza potenzialmente maggiore ma 
    latenza superiore, poco adatto a flussi video. \\
    Metodi tradizionali (background subtraction, ecc.) & Scenario dinamico 
    con camera in movimento e occlusioni: mancanza di \emph{background} stabile. \\
    \bottomrule
  \end{tabularx}
\end{table}


% \fbox{
%   \begin{minipage}{0.9\textwidth}
%     \textbf{Box 3.A -- Scelte progettuali: detector di giocatori}
%     \begin{itemize}
%       \item \textbf{Modello selezionato:} \ac{yolo}v8s. La variante 
%       \emph{small} (s) è stata scelta come punto di partenza per il 
%       suo ottimo equilibrio tra velocità e accuratezza, adatta a 
%       un'analisi quasi in tempo reale (\emph{real-time}) su hardware 
%       consumer.
%       \item \textbf{Alternative scartate:}
%       \begin{itemize}
%         \item \emph{Detector two-stage (ad esempio Faster R-CNN):} 
%         sebbene potenzialmente più accurati, sono stati scartati a 
%         causa della latenza di inferenza significativamente più alta, 
%         che li rende inadatti all'analisi video fluida.
%         \item \emph{Approcci tradizionali (ad esempio \emph{background 
%         subtraction}):} non considerati praticabili a causa della 
%         natura dinamica della scena, con movimenti della telecamera 
%         (anche minimi), cambiamenti di illuminazione e la presenza 
%         costante di giocatori in movimento che rendono impossibile 
%         definire un \emph{background} stabile.
%       \end{itemize}
%     \end{itemize}
%   \end{minipage}
% }
% \todo{valuta se sostituire "real-time" con "tempo reale"} 

% TODO: valutare l'inserimento di una figura esplicativa del razionale di scelta.
%   % TODO: inserisci anche l'immagine, per ora mancante.
% \begin{figure}[htbp]
%   \centering
%   \caption{Razionale dietro la scelta del modello di rilevamento.}
%   \label{box:scelte_detector}
% \end{figure}

\subsection{Addestramento del Modello}
\label{sec:addestramento-modello}
% Descrivere il processo di training: iperparametri utilizzati (learning rate, numero di epoche, batch size),
% l'hardware su cui è stato effettuato l'addestramento, e come è stato suddiviso il dataset (training set,
% validation set, test set).
Il processo di addestramento è stato gestito tramite lo script 
\texttt{\seqsplit{training\_notebooks/basketball\_player\_detection\_training.ipynb}}. 
E' stata adottata una strategia di \emph{transfer learning}, che 
consente di accelerare e migliorare il training sfruttando la conoscenza 
pre-acquisita da un modello addestrato su un dataset di larga scala.

Il punto di partenza è stato un modello \ac{yolo}v8s pre-addestrato 
sul dataset \ac{coco}. I pesi di questo modello sono stati poi affinati 
(\emph{fine-tuning}) utilizzando il dataset specifico 
\texttt{basketball-player-detection-2} descritto nella Sezione~\ref{sec:dataset-utilizzati}. 
Questo approccio consente al modello di adattare le sue \emph{feature} 
gerarchiche, apprese su oggetti generici, al dominio specifico dei 
giocatori di pallacanestro.

Il training è stato configurato con gli iperparametri riportati in 
Tabella~\ref{tab:hyperparams_player_det}. E' stata implementata una 
strategia di \emph{early stopping} con una pazienza di 30 epoche: 
l'addestramento si interrompe automaticamente se la metrica di 
validazione primaria (\ac{map}@0.5:0.95) non migliora sul 
\emph{validation set} per 30 epoche consecutive, salvando il modello 
con le migliori performance e prevenendo l'\emph{overfitting}.

\begin{table}[h]
  \centering
  \caption{Iperparametri usati per l'addestramento del rilevatore 
  dei giocatori (\ac{yolo}v8s).}
  \label{tab:hyperparams_player_det}
  \begin{tabular}{l l}
    \toprule
    \textbf{Parametro} & \textbf{Valore} \\
    \midrule
    Pesi iniziali & \texttt{yolov8s.pt} (pre-addestrato su \ac{coco}) \\
    Epoche & 120 \\
    Pazienza (early stopping) & 30 \\
    Dimensione immagine ($imgsz$) & 960 \\
    Batch & \texttt{-1} (adattivo), con \texttt{accumulate} non usato \\
    Ottimizzatore & \texttt{auto} (impostazione Ultralytics) \\
    $lr0$ (iniziale) & 0.003 \\
    $lrf$ (fattore finale) & 0.2 \\
    Momento & 0.937 \\
    Decadimento pesi & $5\cdot10^{-4}$ \\
    Warmup (epoche) & 3.0 \\
    Multi-scala & \texttt{True} \\
    Mosaic & 1.0 (disattivato nelle ultime 10 epoche, \texttt{close\_mosaic}=10) \\
    MixUp & 0.1 \\
    Flip orizzontale ($fliplr$) & 0.5 \\
    Traslazione / Scala & 0.05 / 0.3 \\
    Rotazioni / Shear / Prospettiva & 0 / 0 / 0 \\
    Max detections ($max\_det$) & 300 \\
    \bottomrule
  \end{tabular}
\end{table}

\noindent L'addestramento è stato eseguito su Google Colab con una 
\ac{gpu} di tipo Tesla T4.
% TODO: potrebbe aver senso introdurre ulteriori dettagli sulla GPU 
% citata, altrimenti direttamente citare il documento colab che ne parla 
% in dettaglio tecnico.

\subsection{Post-processing e Conteggio}
\label{subsec:post-processing-conteggio}

L'output grezzo di un modello di rilevamento (object detection) 
consiste in un numero elevato di potenziali bounding box, molti 
dei quali ridondanti o errati. Per ottenere un output pulito e 
affidabile è essenziale una fase di post-elaborazione (post-processing).

\begin{enumerate}
  \item \textbf{Filtraggio per soglia di confidenza (\emph{confidence 
  thresholding})}: ogni rilevamento prodotto dal modello è associato 
  a un punteggio di confidenza (compreso tra $0$ e $1$) che indica 
  quanto il modello è ``sicuro'' della sua predizione. Viene applicata 
  una soglia (ad esempio $0{,}3$) per scartare tutti i rilevamenti con 
  confidenza inferiore, eliminando gran parte dei falsi positivi.
  \item \textbf{Non-Maximum Suppression (\ac{nms})}: è comune che il 
  modello produca più bounding box, leggermente sfalsati, per lo stesso 
  oggetto. L'algoritmo \ac{nms} risolve questa ambiguità: per ogni gruppo 
  di box sovrapposti che superano una certa soglia di \ac{iou}, conserva 
  solo quello con il punteggio di confidenza più alto e sopprime tutti 
  gli altri. Questo processo, illustrato in Figura~\ref{fig:nms_effect}, 
  è cruciale per evitare doppi conteggi.
\end{enumerate}

\noindent\textbf{Impostazioni adottate.}
Per coerenza in tutti gli esperimenti, salvo diversa indicazione, 
abbiamo utilizzato una soglia di confidenza pari a $0{,}30$ e una 
soglia \ac{iou} per la \ac{nms} pari a $0{,}70$, con NMS sensibile 
alla classe (\texttt{agnostic}=\texttt{False}). Per la Figura~\ref{fig:nms_effect} 
è stato utilizzato il modello generale (\ac{yolo}v8s, 11 classi) con $imgsz=960$; 
l'immagine proviene dal set di validazione e il conteggio è calcolato 
sull'insieme di classi \texttt{player*}.

\noindent\textbf{Regola di conteggio.} Una predizione è conteggiata come 
\emph{giocatore} se la sua classe appartiene a
\[
\mathcal{C}_{\mathrm{player}}=
\left\{
\begin{array}{l}
\texttt{player}, \texttt{player-dribble}, \texttt{player-fall},\\
\texttt{player-jump-shot}, \texttt{player-layup}, \texttt{player-screen}, \texttt{player-shot-block}
\end{array}
\right\}.
\]
Sono escluse \texttt{ball}, \texttt{number}, \texttt{referee}, \texttt{rim}.

Le denominazioni delle classi seguono quelle riportate in 
Figura~\ref{fig:players-actions-classes}.


\begin{table}[h]
  \centering
  \small
  \caption{Mappatura delle classi ai fini del conteggio dei giocatori.}
  \label{tab:classi_conteggio}
  \begin{tabular}{l c}
    \toprule
    \textbf{Classe} & \textbf{Conta come giocatore} \\
    \midrule
    \texttt{player} & Sì \\
    \texttt{player-dribble} & Sì \\
    \texttt{player-fall} & Sì \\
    \texttt{player-jump-shot} & Sì \\
    \texttt{player-layup} & Sì \\
    \texttt{player-screen} & Sì \\
    \texttt{player-shot-block} & Sì \\
    \texttt{ball} & No \\
    \texttt{number} & No \\
    \texttt{referee} & No \\
    \texttt{rim} & No \\
    \bottomrule
  \end{tabular}
\end{table}


\begin{figure}[h]
  \centering
  \includegraphics[width=\textwidth]{figures/03_architettura/nms_effect_final.png}
  \caption{Effetto della \ac{nms} sul rilevamento dei giocatori. 
  A sinistra: predizioni prima della soppressione (conf=0{,}001, \ac{iou}=1{,}00). 
  A destra: risultato dopo \ac{nms} (conf=0{,}30, \ac{iou}=0{,}70, \texttt{agnostic}=\texttt{False}). 
  Il conteggio in basso è calcolato sull'insieme di classi \texttt{player*}.}
  \label{fig:nms_effect}
\end{figure}

\subsection*{Punti chiave}

\begin{itemize}
  \item La scelta di \ac{yolo}v8s bilancia velocità e accuratezza, 
  requisiti fondamentali per l'analisi video.
  \item Il \emph{transfer learning} consente di specializzare un 
  modello potente su un dominio specifico con un effort computazionale ragionevole.
  \item Il post-processing (filtraggio per soglia di confidenza e 
  \ac{nms}) è un passaggio critico per trasformare le predizioni grezze 
  in dati utilizzabili e per un conteggio affidabile.
\end{itemize}

\subsection*{Implicazioni per il Cap.~\ref{chap:risultati-sperimentali}}

\begin{itemize}
  \item Le performance del rilevatore dovranno essere valutate 
  quantitativamente sul test set usando le metriche standard 
  \ac{map}@0.5 e \ac{map}@0.5:0.95.
  \item Un'analisi qualitativa degli errori è necessaria per 
  identificare i limiti del modello, ad esempio in scenari di 
  forte occlusione o con giocatori ai margini del frame.
  \item L'impatto della soglia di confidenza e della soglia 
  \ac{iou} del \ac{nms} sulle metriche di precisione e 
  richiamo (precision-recall) dovrà essere analizzato per 
  individuare il punto di lavoro ottimale.
\end{itemize}
% ++++++++++++++++++++++++++++++++++++++++++++++++++++++++++++++++++++++++++++++++++++++

% ++++++++++++++++++++++++++++++++++++++++++++++++++++++++++++++++++++++++++++++++++++++
\section{Task 2: Rilevamento e Tracciamento della Traiettoria della Palla}
\label{sec:task-2}

Il tracciamento della palla è notoriamente uno dei compiti più 
complessi nell'analisi video sportiva. Le difficoltà principali sono 
legate a tre fattori: (i) dimensioni ridotte (la palla occupa 
pochi pixel), (ii) velocità elevata con frequente \emph{motion blur}, 
(iii) occlusioni ricorrenti dovute ai corpi dei giocatori.
Un tracciamento affidabile è tuttavia indispensabile per abilitare 
l'analisi di eventi di gioco fondamentali (tiri, passaggi, possesso).

\subsection{Approccio implementato}
\label{subsec:approccio-implementato-palla}

Per gestire tali criticità è stato adottato un approccio ibrido 
\emph{tracking-by-detection} a due stadi:

\begin{itemize}
  \item \textbf{Stadio 1 -- Rilevamento (\ac{yolo}v8 dedicato)}: 
  in ogni frame un rilevatore addestrato specificamente sulla classe 
  \texttt{ball} produce una o più ipotesi di localizzazione 
  (\emph{bounding box}). % dataset e procedure in Sez.~\ref{subsec:augmentation}.
  \item \textbf{Stadio 2 -- Stima di stato (Filtro di Kalman)}:
  le osservazioni del rilevatore sono assimilate da un filtro di Kalman
  con modello di moto a velocità costante nel piano immagine 
  (stato: posizione e velocità in $x,y$). Il filtro opera in 
  \emph{predizione} (propagazione con il modello di moto) e 
  \emph{aggiornamento} (correzione con una misura compatibile).
\end{itemize}

In assenza di osservazioni affidabili (occlusioni o outlier) la sola 
predizione mantiene la traiettoria per alcuni frame, garantendo 
\emph{continuità} del tracciamento. 
Per chiarezza espositiva, la Figura~\ref{fig:kalman_sim} mostra, su dati \emph{simulati}, 
come il filtro smussi le misure rumorose e mantenga la traiettoria durante brevi occlusioni.

\begin{figure}[h!]
  \centering
  \includegraphics[width=\textwidth]{figures/03_architettura/kalman_synthetic.png}
  \caption{Esempio \textbf{simulato} dell'effetto del filtro di Kalman sul tracciamento della palla.
  La curva arancione rappresenta la traiettoria generata e i punti arancioni le misure rumorose;
  in blu la stima filtrata; il tratteggio indica la \textbf{predizione in occlusione}. 
  La figura ha scopo illustrativo e \textbf{non} rappresenta un episodio reale del dataset; 
  le valutazioni sperimentali su dati reali sono riportate nel Cap.~\ref{chap:risultati-sperimentali}.}
  \label{fig:kalman_sim}
\end{figure}

\subsection{Dettagli implementativi e mitigazioni}
\label{sec:dettagli-implementativi-palla}

\begin{itemize}
  \item \textbf{Gestione delle occlusioni}. Se il rilevatore non produce 
  una misura compatibile, il sistema procede con la sola predizione. 
  Un contatore di \emph{invisibilità} interrompe la traccia oltre una 
  soglia (10--15 frame), prevenendo la propagazione indefinita.
  \item \textbf{Disambiguazione/outlier}. È applicato un \emph{gating} 
  statistico (distanza di Mahalanobis) tra misura osservata e posizione 
  predetta: le osservazioni troppo distanti sono scartate come falsi 
  positivi (ad es. riflessi, teste, luci).
  \item \textbf{Robustezza al \emph{motion blur}}. La resistenza al 
  blur è perseguita in fase di addestramento tramite augmentation 
  fotometrica e di sfocatura (Sez.~\ref{subsec:augmentation}); il 
  filtro di Kalman contribuisce smussando la traiettoria.
\end{itemize}

\begin{riskbox}{Rischi \& Mitigazioni: tracciamento della palla}
  \label{box:rischi-mitigazioni-palla}
  \begin{itemize}
    \item \textbf{Rischi}: occlusioni prolungate (\emph{track} persa), 
    motion blur estremo (mancato rilevamento), falsi positivi (oggetti 
    rotondi).
    \item \textbf{Mitigazioni}: predizione del filtro durante le occlusioni; 
    augmentation con \emph{blur} in addestramento; \emph{gating} di Kalman 
    per scartare misure incoerenti.
  \end{itemize}
\end{riskbox}

\subsection{Dalla traiettoria agli eventi: rilevamento del tiro}
\label{sec:stima-traiettoria-eventi}

La sequenza delle posizioni stimate dal filtro costituisce una 
\textbf{traiettoria liscia} su cui si innesta il rilevamento di eventi. 
Nel sistema implementato il modulo \texttt{ShotDetector} adotta 
regole deterministiche:

\begin{enumerate}
  \item \textbf{Trigger di inizio tiro}: la velocità verticale stimata 
  diventa positiva (\emph{lancio verso l'alto}) e la palla non è in 
  possesso di un giocatore (il centro della palla non è contenuto nel box 
  del giocatore più vicino).
  \item \textbf{Conferma}: la palla supera in altezza la quota del 
  canestro nel frame (proiezione 2D della quota reale).
  \item \textbf{Esito} (\emph{made/missed}): si valuta la distanza 
  minima tra traiettoria e centro del canestro; se inferiore a un raggio 
  $r=30$~px, il tiro è considerato segnato.
\end{enumerate}

Come verifica aggiuntiva, in assenza di rumore eccessivo, i campioni del 
tratto in volo possono essere approssimati da una \emph{parabola} in 2D, 
caratteristica del moto balistico.

\paragraph{Classificazione 2/3 punti.}
La tipologia del tiro è determinata dalla posizione del tiratore in 
\emph{vista tattica} (metri); se la distanza minima dal canestro più 
vicino è $\geq 6{,}75$~m (standard della \ac{fiba}) il tiro è classificato da tre 
punti, altrimenti da due. In assenza temporanea di coordinate affidabili, 
è disponibile un meccanismo di \emph{ripiego} in 2D.

\subsection*{Parametri principali e riproducibilità}
\label{subsec:params-palla}
I parametri di addestramento del modello palla sono elencati in 
Sez.~\ref{subsec:augmentation} (configurazione \emph{ball-only}).
Qui si riportano le impostazioni di inferenza e del filtro di Kalman 
utilizzate negli esperimenti, per facilitarne la replica.

\begin{table}[h]
  \centering
  \small
  \caption{Parametri operativi per il tracciamento della palla (inferenza e filtro).}
  \label{tab:params_ball_tracking}
  \begin{tabular}{l l}
    \toprule
    \textbf{Parametro} & \textbf{Valore tipico} \\
    \midrule
    Soglia confidenza rilevatore & $0{,}25$ \\
    Gating (chi-quadro, 2D) & $5{,}991$ (95\%) -- $9{,}21$ (99\%) \\
    Rumore misura base ($\sigma_M$) & $5$~px (adattivo in base al box) \\
    Rumore processo ($\sigma_A$) & $120$--$200$~px/s$^2$ \\
    Passo temporale ($\Delta t$) & $1/\mathrm{FPS}$ del video \\
    Invisibilità massima & 10--15 frame \\
    \bottomrule
  \end{tabular}
\end{table}

\subsection*{Punti chiave}
\begin{itemize}
  \item L'approccio ibrido \ac{yolo}v8 + Kalman fornisce un tracciamento 
  robusto, mantenendo la continuità nelle brevi occlusioni.
  \item Il \emph{gating} statistico limita i falsi positivi; la traiettoria 
  stimata abilita regole fisiche e geometriche per il riconoscimento dei tiri.
  \item La riproducibilità è garantita dal riuso degli stessi parametri 
  operativi (Tab.~\ref{tab:params_ball_tracking}) e dalla configurazione 
  di addestramento documentata in Sez.~\ref{subsec:augmentation}.
\end{itemize}

\subsection*{Implicazioni per il Capitolo~\ref{chap:risultati-sperimentali}}
\begin{itemize}
  \item Valutazione del rilevatore (precisione/richiamo) e dell'errore 
  medio (\ac{mse}) sulla traiettoria rispetto a un sottoinsieme annotato.
  \item Valutazione delm \texttt{ShotDetector} (tiri corretti, mancati, falsi).
  \item Analisi dei fallimenti (occlusioni lunghe, passaggi rapidi scambiati 
  per tiri) e sensibilità ai parametri di gating.
\end{itemize}
% ++++++++++++++++++++++++++++++++++++++++++++++++++++++++++++++++++++++++++++++++++++++


% ++++++++++++++++++++++++++++++++++++++++++++++++++++++++++++++++++++++++++++++++++++++
\section{Task 3: Rilevamento dei Punti Chiave del Campo e Stima dell'Omografia}
\label{sec:task-3}

L'obiettivo di questo modulo è ricostruire in ogni frame la geometria del campo
(perimetro, area dei tre secondi, linea da tre punti, asse di metà campo, ecc.)
a partire da un insieme di \textbf{punti chiave} annotati.
La stima della trasformazione proiettiva (\textbf{omografia}) che mappa il piano
immagine sul piano del campo consente di proiettare le traiettorie di palla e
giocatori in una \emph{vista tattica} metrica (top-down), abilitando analisi
spazio-temporali coerenti.

\subsection{Scelta del modello e del formato}
\label{subsec:scelta-modello-keypoints}

Per il rilevamento dei riferimenti geometrici è stato impiegato \textbf{\ac{yolo}v8-pose}
(addestramento per punti chiave) con \textbf{41 keypoint} e uno \emph{scheletro} fisso
(Fig.~\ref{fig:keypoints-template}), definito in Roboflow.
La scelta è motivata da:
\begin{itemize}
  \item \textbf{Coerenza strutturale}: l'uso di un template con \ac{id} univoco per punto
  riduce ambiguità e consente un accoppiamento diretto con il modello geometrico del campo.
  \item \textbf{Robustezza}: \ac{yolo}v8-pose permette la predizione simultanea di più punti anche in presenza di parziali occlusioni.
  \item \textbf{Compatibilità}: lo stesso ecosistema software del Task~\ref{sec:task-1} semplifica flusso e riproducibilità.
\end{itemize}

\paragraph{Dataset.}
Si utilizza il dataset \emph{Basketball Court Detection} arricchito (Sez.~\ref{sec:dataset-utilizzati}),
per un totale di 2525 immagini annotate con lo \emph{scheletro} a 41 punti.

\subsection{Addestramento: impostazioni rilevanti}
\label{subsec:train-keypoints}

L'addestramento è stato condotto in coerenza con il principio di preservare la
geometria del campo: si evitano trasformazioni che alterino le relazioni
proiettive.

\begin{table}[h]
  \centering
  \caption{Iperparametri usati per il rilevamento dei punti chiave del campo (\ac{yolo}v8-pose).}
  \label{tab:hyperparams_keypoints}
  \begin{tabular}{l l}
    \toprule
    \textbf{Parametro} & \textbf{Valore} \\
    \midrule
    Pesi iniziali & \texttt{yolov8s-pose.pt} (pre-addestrato) \\
    Epoche / Pazienza & 120 / 30 \\
    Dimensione immagine ($imgsz$) & 960 \\
    Batch & \texttt{-1} (adattivo) \\
    Ottimizzatore / LR ($lr0$, $lrf$) & \texttt{auto} / 0.003, 0.2 \\
    Momentum / Weight decay & 0.937 / $5\cdot10^{-4}$ \\
    Warmup (epoche) & 3.0 \\
    Multi-scala & \texttt{False} \\
    Mosaic / MixUp / Copy-Paste & 0.0 / 0.0 / 0.0 \\
    Flip orizzontale ($fliplr$) & 0.0 \\
    HSV (h, s, v) & 0.015, 0.5, 0.3 \\
    Traslazione / Scala & 0.05 / 0.2 \\
    Rotazioni / Shear / Prospettiva & 0 / 0 / 0 \\
    \bottomrule
  \end{tabular}
\end{table}

\noindent
\textit{Razionale.} \emph{Mosaic}, rotazioni e prospettive artificiali sono disabilitate
per non introdurre configurazioni non fisiche del campo. Il flip orizzontale è disattivato
perché altererebbe la semantica sinistra/destra dei punti (gli \ac{id} restano invariati).
Variazioni fotometriche moderate migliorano la robustezza a illuminazione indoor.

\subsection{Dalle predizioni all'omografia: pipeline}
\label{subsec:homography-pipeline}

Dato un frame $t$, il modello produce un sottoinsieme $V_t$ dei 41 punti con
coordinate immagine $\mathbf{p}_i=(x_i,y_i)$ e \ac{id} noto.
Sia $\mathbf{P}_i=(X_i,Y_i)$ il punto corrispondente del template metrico
(28$\times$15 m, standard della \ac{fiba}). La stima dell'omografia $H_t$ segue:

\begin{enumerate}
  \item \textbf{Selezione corrispondenze valide}:
  si usano le coppie $(\mathbf{p}_i,\mathbf{P}_i)$ con confidenza sopra soglia
  e geometria non degenere (almeno 4 punti non collineari).
  \item \textbf{Stima robusta}:
  omografia $H_t$ tramite \ac{dlt} con \ac{ransac}
  (funzione \texttt{findHomography}, soglia di reproiezione $\epsilon$ in px,
  confidenza 0.995), ottenendo inlier/outlier.
  \item \textbf{Rifinitura}:
  ottimizzazione su tutti gli inlier (\ac{lmeds} o minimizzazione non lineare dell'errore di reproiezione).
  \item \textbf{Controllo qualità}:
  accettazione se \#inlier $\geq N_{\min}$ e \ac{rmse}$_\mathrm{repr}\leq\tau$; in caso contrario,
  si mantiene $H_{t-1}$ (\emph{hold}) o si interpola (filtro \ac{ema} sui parametri).
\end{enumerate}

La trasformazione di un punto immagine $\mathbf{p}$ in coordinate metriche
segue $\tilde{\mathbf{P}} \sim H_t\,\tilde{\mathbf{p}}$ (coordinate omogenee).
Inversamente, si può \emph{warpare} l'immagine per ottenere la vista top-down.

\subsection{Stabilizzazione temporale e casi limite}
\label{subsec:stabilization}

Per ridurre jitter e piccoli salti dovuti alla variazione del set di inlier:
\begin{itemize}
  \item \textbf{Smussamento} di $H_t$ con filtro \ac{ema} ($\alpha\in[0.1,0.3]$) o,
  in alternativa, smussamento direttamente sulle traiettorie proiettate.
  \item \textbf{Coerenza di orientamento}:
  l'\emph{ID} dei punti (es. vertici delle aree) determina in modo univoco verso
  e orientazione; in caso di ambiguità si risolve imponendo vincoli (metà campo,
  lato panchine) oppure confrontando la direzione media delle traiettorie.
  \item \textbf{Fallback}:
  se \#inlier $<N_{\min}$ o \ac{rmse} sfora $\tau$, si mantiene $H_{t-1}$ per un numero
  limitato di frame; oltre tale soglia si invalida la vista tattica per quel tratto.
\end{itemize}

\subsection{Parametri operativi e riproducibilità}
\label{subsec:params-homography}

\begin{table}[h]
  \centering
  \small
  \caption{Parametri principali per la stima dell'omografia e la stabilizzazione.}
  \label{tab:params_homography}
  \begin{tabular}{l l}
    \toprule
    \textbf{Parametro} & \textbf{Valore tipico} \\
    \midrule
    Soglia confidenza keypoint & 0.25 \\
    Min corrispondenze non collineari ($N_{\min}$) & 8--12 \\
    Metodo \ac{ransac} / confidenza & \ac{ransac} / 0.995 \\
    Soglia reproiezione ($\epsilon$) & 3--5 px \\
    \ac{rmse} massimo accettato ($\tau$) & 4--6 px \\
    Smussamento (\ac{ema} su $H$) & $\alpha=0.2$ \\
    Coordinate template & \ac{fiba} 28$\times$15 m (metri) \\
    \bottomrule
  \end{tabular}
\end{table}

\subsection{Output e integrazione nella pipeline}
\label{subsec:output-integrazione}

Il modulo restituisce per ciascun frame:
\begin{enumerate}
  \item l'omografia $H_t$ (e la sua qualità: \#inlier, \ac{rmse});
  \item l'insieme dei keypoint visibili con confidenza;
  \item la \textbf{vista tattica}: proiezione in metri delle posizioni di palla e giocatori
  (Task~\ref{sec:task-1} e Task~\ref{sec:task-2}) tramite $H_t$; opzionalmente,
  l'immagine warppata (top-down) per verifiche qualitative.
\end{enumerate}

\begin{riskbox}{Rischi \& Mitigazioni: geometria del campo}
  \label{box:rischi-mitigazioni-campo}
  \begin{itemize}
    \item \textbf{Rischi}: poche linee visibili, forti occlusioni, riflessi; distorsione ottica
    non modellata; camera con forte rollio/pitch variabile.
    \item \textbf{Mitigazioni}: soglie conservative su inlier/\ac{rmse}; smussamento temporale;
    esclusione dei keypoint meno affidabili; eventuale pre-correzione di distorsione.
  \end{itemize}
\end{riskbox}

\subsection*{Punti chiave}
\begin{itemize}
  \item L'uso di \ac{yolo}v8-pose con \emph{ID} semantici permette un accoppiamento diretto
  con il template metrico e una stima di omografia robusta.
  \item La stabilizzazione temporale di $H_t$ riduce jitter e migliora la consistenza della
  vista tattica, facilitando misure e confronti nel tempo.
  \item Parametri e soglie (Tab.~\ref{tab:hyperparams_keypoints}-\ref{tab:params_homography})
  rendono la procedura pienamente riproducibile.
\end{itemize}

\subsection*{Implicazioni per il Capitolo~\ref{chap:risultati-sperimentali}}
\begin{itemize}
  \item \textbf{Accuratezza keypoint}: PCK@$\delta$ (in px o \% diagonale) e media \ac{rmse}
  di reproiezione sugli inlier.
  \item \textbf{Stabilità omografia}: varianza dei parametri di $H_t$ o dell'errore di
  reproiezione lungo il tempo; tasso di \emph{fallback}.
  \item \textbf{Qualità vista tattica}: errore medio sulla posizione di palla/giocatori
  in metri (dove disponibile ground truth) e coerenza geometrica con le dimensioni FIBA.
\end{itemize}
% ++++++++++++++++++++++++++++++++++++++++++++++++++++++++++++++++++++++++++++++++++++++


% ++++++++++++++++++++++++++++++++++++++++++++++++++++++++++++++++++++++++++++++++++++++
\section{Integrazione, Ottimizzazioni e Deployment}
\label{sec:integrazione-ottimizzazioni}

Lo sviluppo dei singoli moduli è solo una parte del lavoro: per ottenere un sistema
utilizzabile i componenti devono essere integrati in una pipeline coerente, ottimizzati
per funzionare in modo efficiente e confezionati per un'esecuzione riproducibile.

\subsection{Orchestrazione della pipeline}
\label{subsec:orchestrazione-pipeline}

L'integrazione dei moduli è gestita da uno \emph{script} principale (\texttt{main.py}) che
inizializza: (i) i \emph{detector} (giocatori, palla, punti chiave del campo),
(ii) i \emph{tracker} (filtro di Kalman e associazione), (iii) i calcolatori di metriche,
(iv) i visualizzatori (\texttt{drawers}). Il ciclo di elaborazione in modalità \emph{batch} è:

\begin{enumerate}
  \item apertura del video d'ingresso e iterazione sui frame;
  \item per ogni frame, invocazione in sequenza dei moduli (Sez.~\ref{sec:architettura-sistema}):
  punti chiave del campo $\rightarrow$ omografia (Task~\ref{sec:task-3}) $\rightarrow$
  rilevamento giocatori (Task~\ref{sec:task-1}) e palla (Task~\ref{sec:task-2}) $\rightarrow$
  tracciamento/associazione $\rightarrow$ aggiornamento eventi;
  \item scrittura del frame annotato nel video di uscita.
\end{enumerate}

Questo approccio è ideale per l'analisi post-partita. Un'estensione al \emph{tempo reale}
richiederebbe sorgenti \emph{stream} e una pipeline con precaricamento, mini-lotti (\emph{mini-batch})
e parallelizzazione per ridurre la latenza end-to-end.

\subsection{Profilazione e ottimizzazione delle performance}
\label{sec:profilazione-ottimizzazione}

La profilazione mostra che l'inferenza dei modelli di \emph{deep learning} domina il tempo
di calcolo, seguita dai trasferimenti \ac{cpu}$\leftrightarrow$\ac{gpu} e dall'\ac{io} video.

\paragraph{Strategie pratiche.}
\begin{itemize}
  \item \textbf{FP16} (mezza precisione) per tutti i modelli compatibili: in genere aumenta gli \ac{fps}
  con impatto trascurabile sull'accuratezza.
  \item \textbf{Esportazione ONNX / TensorRT}: fusione di layer e kernel ottimizzati per la \ac{gpu},
  particolarmente efficace per i detector invocati a ogni frame.
  \item \textbf{Mini-lotti in inferenza} (ad es. 4-8) quando non è richiesta una dipendenza temporale stretta; i risultati vengono poi riordinati.
  \item \textbf{Trasferimenti asincroni} e memoria \emph{pinned} per ridurre l'overhead
  \ac{cpu}$\to$\ac{gpu}; \emph{warm-up} iniziale per stabilizzare le latenze.
  \item \textbf{Decodifica/ricodifica accelerate} (NVDEC/NVENC) o \texttt{ffmpeg} con
  \texttt{hevc\_nvenc}; \emph{prebuffer} in lettura. L'uso di immagini rettangolari (\texttt{rect})
  riduce il padding in inferenza.
\end{itemize}

\paragraph{Setup di misura.}
Le misure di latenza e throughput sono state raccolte \textbf{direttamente su Google Colab} utilizzando una
\textbf{GPU NVIDIA Tesla T4} (CUDA~12.6, VRAM~$\sim$15{,}8~GB). Il video sorgente è 1280$\times$720@30~fps.
I modelli impiegati sono YOLOv8s (giocatori, $imgsz{=}960$, FP16) e YOLOv8m (palla, $imgsz{=}1280$, FP16).

\begin{table}[h]
  \centering
  \caption{Profilazione per stadio su \textbf{Google Colab} (GPU \textbf{NVIDIA Tesla T4}, CUDA~12.6) con input 1280$\times$720.
  YOLOv8s (giocatori, $imgsz{=}960$, FP16) e YOLOv8m (palla, $imgsz{=}1280$, FP16).}
  \label{tab:latenze_profiling_t4}
  \begin{tabular}{lrr}
    \toprule
    \textbf{Stadio} & \textbf{Media [ms]} & \textbf{Quota [\%]} \\
    \midrule
    Decodifica video & 1.88 & 4.5 \\
    Inferenza giocatori (end-to-end) & 18.29 & 44.2 \\
    Inferenza palla (end-to-end) & 21.23 & 51.3 \\
    \midrule
    \textbf{Totale per frame} & \textbf{41.39} & \textbf{100.0} \\
    \bottomrule
  \end{tabular}
  \vspace{0.3em}

  \small Throughput misurato: \textbf{24.14 FPS} su 305 frame (profilazione eseguita su runtime Google Colab).
\end{table}

\noindent\emph{Nota.} I valori riportati riflettono il runtime Google Colab con GPU Tesla T4 e possono differire da esecuzioni su hardware locale.

\begin{riskbox}{Criticità \& mitigazioni nel deployment}
  \label{box:rischi-mitigazioni-deploy}
  \begin{itemize}
    \item \textbf{Degrado di accuratezza in FP16/INT8}: verificare che la variazione di \ac{map} sia $<0{,}5$ p.p.; in caso contrario,
    usare FP32 per i soli modelli sensibili.
    \item \textbf{Collo di bottiglia \ac{io}}: abilitare NVDEC/NVENC e \emph{prebuffer}; evitare ricampionamenti costosi nel loop.
    \item \textbf{Saturazione della VRAM}: ridurre $imgsz$ o il batch, abilitare mini-lotti; spostare la grafica di overlay su \ac{cpu} se necessario.
  \end{itemize}
\end{riskbox}

\subsection{Packaging e requisiti di esecuzione}
\label{sec:packaging-deployment}

Il progetto è containerizzato con Docker per assicurare replicabilità:
il \texttt{Dockerfile} crea un'immagine con sistema, dipendenze di sistema
(OpenCV, \texttt{ffmpeg}, \texttt{TensorRT} opzionale), ambiente Python e librerie
(\texttt{requirements.txt}). Requisiti consigliati:
\begin{itemize}
  \item \ac{gpu} NVIDIA con supporto \ac{cuda} e $\geq$8\,GB di \ac{vram};
  \item 16\,GB di \ac{ram};
  \item Linux/Windows con Docker e NVIDIA Container Toolkit.
\end{itemize}

Il throughput effettivo dipende dall'hardware e dalla configurazione; i valori misurati su Colab (Tab.~\ref{tab:latenze_profiling_t4})
costituiscono il riferimento per questa tesi.

\subsection{Manifest di riproducibilità}
\label{subsec:manifest-reprod}

Per replicare gli esperimenti:
\begin{itemize}
  \item \textbf{Versioni}: \texttt{ultralytics}\,/\,\texttt{torch}\,/\,\texttt{cuda} (riportate nell'Appendice tecnica);
  \item \textbf{Dataset}: slug/ID Roboflow e versione di esportazione (Sez.~\ref{sec:dataset-utilizzati} e Tab.~\ref{tab:dataset_split});
  \item \textbf{Modelli}: pesi ``\emph{best}'' per i tre task, con \ac{yolo}v8s (giocatori), \ac{yolo}v8m (palla), \ac{yolo}v8-pose (punti chiave);
  \item \textbf{Parametri}: soglie e impostazioni operative come riportate nelle tabelle dei Task~\ref{sec:task-1}-\ref{sec:task-3};
  \item \textbf{Seed}: inizializzazione fissata, dove applicabile, per ridurre la variabilità in addestramento.
\end{itemize}

\subsection*{Punti chiave}
\begin{itemize}
  \item Pipeline \emph{batch} orchestrata da uno \emph{script} unico: semplice, riproducibile, adatta all'analisi offline.
  \item FP16, ONNX/TensorRT, mini-lotti e codec accelerati sono i principali moltiplicatori di \emph{throughput}.
  \item La containerizzazione con Docker garantisce portabilità e coerenza dell'ambiente.
\end{itemize}

\subsection*{Implicazioni per il Cap.~\ref{chap:risultati-sperimentali}}
\begin{itemize}
  \item Riportare latenza end-to-end e \ac{fps} del sistema integrato (Tab.~\ref{tab:latenze_profiling_t4}) oltre alle metriche di accuratezza.
  \item Confrontare PyTorch nativo vs ONNX/TensorRT (e FP32 vs FP16) per quantificare i guadagni.
  \item Discutere il \emph{trade-off} tra ricchezza delle analisi prodotte e costo computazionale.
\end{itemize}
% ++++++++++++++++++++++++++++++++++++++++++++++++++++++++++++++++++++++++++++++++++++++



%----------------------------------------------------------------------------------------
\chapter{Risultati Sperimentali e Valutazione (10-20 pagine)}
\label{chap:risultati-sperimentali}
%----------------------------------------------------------------------------------------

Questo capitolo presenta i risultati quantitativi e qualitativi ottenuti dall'applicazione della metodologia 
descritta nel Capitolo~\ref{chap:metodologia-sviluppo}. L'obiettivo è validare empiricamente le scelte progettuali, 
misurare le performance di ogni componente chiave del sistema e confrontare i risultati finali con i criteri di 
successo predefiniti.

L'analisi è strutturata per seguire il piano sperimentale delineato nella Sezione \ref{sec:controllo-qualita-esperimenti}. 
Si inizierà con la definizione formale delle metriche di valutazione, per poi passare all'analisi dei risultati 
del rilevamento dei giocatori e del tracciamento della palla, includendo gli esiti degli studi di ablazione. 
Infine, i risultati verranno discussi criticamente, evidenziando i punti di forza e i limiti del sistema sviluppato.

% ++++++++++++++++++++++++++++++++++++++++++++++++++++++++++++++++++++++++++++++++++++++
% Richiesta:
% Metriche di Valutazione: Definire le metriche usate per misurare le performance.
% \begin{itemize}
%   \item \textbf{Per il Rilevamento Giocatori:} Errore medio percentuale rispetto al ground truth?.
%   \item \textbf{Per la Traiettoria della Palla:} Errore Quadratico Medio (MSE) tra le coordinate
%   predette e quelle del ground truth?
% \end{itemize}
\section{Metriche di Valutazione:}
\label{sec:metriche-valutazione}
Per garantire una valutazione oggettiva e standardizzata, sono state utilizzate diverse metriche, ciascuna adatta 
a misurare un aspetto specifico del sistema.

\subsection{Metriche per il Rilevamento di Oggetti (Detection)}
\label{sec:metriche-detection}
Il task di rilevamento (giocatori, palla) è stato valutato utilizzando metriche standard basate sul concetto di 
\ac{iou}, che misura il grado di sovrapposizione tra un bounding box predetto e uno di ground truth.

\paragraph{Precision (Precisione).} proporzione di predizioni positive corrette sul totale delle predizioni positive; 
valori elevati indicano pochi falsi positivi.
\begin{equation}
  \label{eq:precision}
  {Precision} = \frac{TP}{TP + FP}\footnotemark
\end{equation}
\footnotetext{TP = True Positive (vero positivo); FP = False Positive (falso positivo).}

\paragraph{Recall (Copertura).} proporzione di positivi reali correttamente rilevati sul totale dei positivi reali; 
valori elevati indicano pochi falsi negativi.
\begin{equation}
  \label{eq:recall}
  {Recall} = \frac{TP}{TP + FN}\footnotemark
\end{equation}
\footnotetext{TP = True Positive (vero positivo); FN = False Negative (falso negativo).}

\paragraph{\acf{map}.} metrica di riferimento per l'object detection. 
Per ogni classe $(i)$, l'\ac{ap} è l'area sotto la curva precision-recall; 
la \ac{map} è la media delle $(AP\_i)$ su tutte le $(n)$ classi.
\begin{itemize}
  \item \ac{map}@0.50: calcolata con soglia \ac{iou} fissa a 0.50; valuta la capacità di localizzare 
  correttamente gli oggetti con tolleranza moderata.
  \item \ac{map}@0.50:0.95: media dell'\ac{ap} su dieci soglie \ac{iou} da 0.50 a 0.95 con passo 0.05 
  (definizione in stile \ac{coco}); valuta la precisione di localizzazione più stringente.
\end{itemize}
\begin{equation}
  \label{eq:map}
  {mAP} = \frac{1}{n} \sum_{i=1}^{n} AP\_i \footnotemark
\end{equation}
\footnotetext{$n$ = numero di classi; $AP\_i$ = Average Precision della classe $i$.}

\subsection{Metriche per la stima di traiettorie}
\label{sec:metriche-traiettorie}
Per valutare l'accuratezza della traiettoria della palla utilizziamo la \ac{rmse}, calcolata sulla distanza euclidea
tra le coordinate predette e quelle di ground truth lungo l'intera traiettoria. E' espressa in pixel e fornisce una misura
intuitiva dell'errore medio di localizzazione.

\paragraph{\ac{rmse}.}
\begin{equation}
  \label{eq:rmse_2d}
  {RMSE} = \sqrt{\frac{1}{n}\sum_{i=1}^{n}\Big[(x_i - \hat{x}_i)^2 + (y_i - \hat{y}_i)^2\Big]}\,\footnotemark
\end{equation}
\footnotetext{$n$ = numero di campioni (frame) della traiettoria; $(x_i, y_i)$ = coordinate di ground truth al 
frame $i$; $(\hat{x}_i, \hat{y}_i)$ = coordinate stimate al frame $i$.}

\subsection{Metriche di performance del sistema}
\paragraph{\acf{fps}.} Misura il \emph{throughput} del sistema, ovvero il numero di frame elaborati al secondo; 
è l'indicatore chiave della velocità ed efficienza computazionale.
\begin{equation}
  \label{eq:fps}
  {FPS} = \frac{n}{\Delta t}\,\footnotemark
\end{equation}
\footnotetext{$n$ = numero di frame effettivamente processati; $\Delta t$ = intervallo temporale misurato 
(tempo di elaborazione netto sul periodo considerato, tipicamente esclusi warm-up e \ac{io}, salvo diversa indicazione).}

\noindent In modalità sequenziale (batch $(=1)$), la latenza media per frame è $(\bar{t}_{\text{frame}}=\tfrac{1}{n}\sum_i t_i)$ 
e vale approssimativamente $({FPS}\approx 1/\bar{t}_{\text{frame}})$ (oppure $({ms/frame}\approx 1000/\ac{fps})$). 
In presenza di pipeline o parallelismo, throughput e latenza non coincidono necessariamente: è quindi buona prassi 
riportare anche la latenza end-to-end.
% ++++++++++++++++++++++++++++++++++++++++++++++++++++++++++++++++++++++++++++++++++++++

% ++++++++++++++++++++++++++++++++++++++++++++++++++++++++++++++++++++++++++++++++++++++
% Richiesta:
% Presentare i risultati in modo chiaro usando tabelle e grafici. Confrontare le performance ottenute sul
% dataset pubblico e sui video registrati. Discutere eventuali differenze. Includere esempi visivi 
% (frame con i bounding box disegnati dal modello).
\section{Risultati del Rilevamento Giocatori}
\label{sec:risultati-rilevamento-giocatori}
Questa sezione presenta le performance del modello \ac{yolo}v8m per il rilevamento dei giocatori, 
con studi di ablazione a supporto delle scelte metodologiche, e un confronto fra il dataset 
pubblico e i video registrati.

\subsection{Performance del Modello Finale}
Il modello finale è stato addestrato con data augmentation \texttt{aug3x} e una risoluzione di input di $640\times640$ pixel. 
La \autoref{tab:player_detector_final_results} riassume le metriche sul test set.
% fa riferimento alla tabella 4.1

\begin{table}[t]
  \centering
  \caption{Performance del modello \ac{yolo}v8m finale sul test set. 
  Le metriche superano i criteri di successo (\ac{map}@0.5 $\geq$ 92\%).}
  \label{tab:player_detector_final_results}
  \begin{tabular}{l S[table-format=1.3] S[table-format=1.3] S[table-format=1.3] S[table-format=1.3] S[table-format=2.1]}
    \toprule
    & {\textbf{Precision}} & {\textbf{Recall}} & {\textbf{\ac{map}@0.5}} & {\textbf{\ac{map}@0.5:0.95}} & {\textbf{\ac{fps}}} \\
    \midrule
    \textbf{Finale (aug3x, 640)} 
      & 0.96 & 0.94 & 0.955 & 0.825 & 45.3 \\
    \bottomrule
  \end{tabular}
\end{table}

% TODO: Decommenta il codice LaTeX sottostante per includere la tabella reale.
\begin{table}[h]
  \centering
  % \begin{tabular}{l c c c c} ... \end{tabular}
  \framebox[1.0\textwidth][c]{Tabella con i risultati del modello finale di rilevamento giocatori}
  \caption{Performance del modello \ac{yolo}v8m finale sul test set. Le metriche superano i criteri di successo definiti 
  (\ac{map}@0.5 $\geq$ 92\%), confermando l'elevata accuratezza del detector.}
  \label{tab:player_detector_final_results}
\end{table}
% TODO: sostituire il placeholder con una tabella reale (tabular) e assicurare coerenza con il testo.

Curva Precision-Recall (opzionale, o includi un'immagine esterna)

% TODO: Decommenta codice LaTeX per includere la curva PR, oppure inserisci un'immagine esterna.}
\begin{figure}[t]
  \centering
  % Opzione A: genera il grafico da CSV (results/pr_curve.csv con colonne: recall,precision)
  % \begin{tikzpicture}
  %   \begin{axis}[
  %     width=0.7\linewidth, height=6cm,
  %     xlabel=Recall, ylabel=Precision,
  %     xmin=0, xmax=1, ymin=0, ymax=1,
  %     grid=both, legend style={at={(0.97,0.2)},anchor=east}
  %   ]
  %   \addplot+[mark=none] table[x=recall,y=precision,col sep=comma]{results/pr_curve.csv};
  %   \addlegendentry{PR curve (test set)}
  %   \end{axis}
  % \end{tikzpicture}
  % Opzione B: includi un PDF/PNG già esportato
  % \includegraphics[width=0.7\linewidth]{figures/pr_curve_final.pdf}
  % TODO: Curva Precision-Recall del modello finale sul test set. 
  % Generare il grafico da results/pr_curve.csv e esportarlo come PDF.}
  \caption{Curva Precision-Recall del modello finale sul test set.}
  \label{fig:pr_curve_final}
\end{figure}

I risultati indicano alta \emph{Recall} (rileva la maggior parte dei giocatori presenti) e alta \emph{Precision} 
(pochi falsi positivi), con \ac{map}@0.5 $>\,$0.95.

\subsection{Studio di Ablazione 1: Impatto della Data Augmentation}
Per quantificare il beneficio della data augmentation, è stato addestrato un modello \texttt{Baseline} senza 
le trasformazioni della Sezione~\autoref{sec:augmentation}.
La \autoref{tab:ablation_augmentation} confronta le performance.
% (Tabella 4.2) confronta le performance dei due modelli.

\begin{table}[t]
  \centering
  \caption{Confronto tra \texttt{Baseline} (senza augmentation) e \texttt{aug3x}. 
  L'augmentation migliora tutte le metriche.}
  \label{tab:ablation_augmentation}
  \begin{tabular}{l S[table-format=1.3] S[table-format=1.3] S[table-format=1.3] 
    S[table-format=1.3] S[table-format=2.1]}
    \toprule
    & {\textbf{Precision}} & {\textbf{Recall}} & {\textbf{\ac{map}@0.5}} & 
    {\textbf{\ac{map}@0.5:0.95}} & {\textbf{\ac{fps}}} \\
    \midrule
    \texttt{Baseline} & 0.89 & 0.82 & 0.878 & 0.695 & 45.3 \\
    \texttt{aug3x}    & 0.96 & 0.94 & 0.955 & 0.825 & 45.3 \\
    \bottomrule
  \end{tabular}
\end{table}
\missingfigure{inserisci immagine comparativa tra i due modelli}

\begin{figure}[t]
  \centering
  % Grafico a barre raggruppate (Baseline vs aug3x)
  \begin{tikzpicture}
    \begin{axis}[
      width=0.85\linewidth, height=6cm,
      ybar, bar width=12pt,
      symbolic x coords={Precision,Recall,\ac{map}@0.5,\ac{map}@0.5:0.95},
      xtick=data, ymin=0, ymax=1,
      legend style={at={(0.5,1.05)}, anchor=south, legend columns=-1},
      ylabel={Valore},
      grid=both
    ]
      \addplot coordinates {(Precision,0.89) (Recall,0.82) (\ac{map}@0.5,0.878) 
      (\ac{map}@0.5:0.95,0.695)};
      \addplot coordinates {(Precision,0.96) (Recall,0.94) (\ac{map}@0.5,0.955) 
      (\ac{map}@0.5:0.95,0.825)};
      \legend{Baseline, aug3x}
    \end{axis}
  \end{tikzpicture}
  \caption{Confronto metrico fra \texttt{Baseline} e \texttt{aug3x}.}
  \label{fig:bar_ablation}
\end{figure}
\missingfigure{testo che spiega il grafico}

Si osserva un incremento consistente in \ac{map}@0.5 ($\approx +8$ punti) e \ac{map}@0.5:0.95 ($\approx +13$ punti), 
coerente con l'ipotesi che l'augmentation (variazioni di luminosità, blur, occlusioni) migliori robustezza e accuratezza.

\subsection{Confronto: Dataset Pubblico vs Video Registrati}
Per valutare il \emph{domain gap}, riportiamo i risultati su due insiemi: il dataset pubblico annotato e i video registrati in contesto reale. 
La \autoref{tab:domain_gap} e la \autoref{fig:domain_gap_bar} mostrano differenze e trend.

\begin{table}[t]
  \centering
  \caption{Confronto delle performance tra dataset pubblico e video registrati (stesso modello).}
  \label{tab:domain_gap}
  \begin{tabular}{l S[table-format=1.3] S[table-format=1.3] S[table-format=1.3] S[table-format=1.3]}
    \toprule
    & {\textbf{Precision}} & {\textbf{Recall}} & {\textbf{\ac{map}@0.5}} & {\textbf{\ac{map}@0.5:0.95}} \\
    \midrule
    Pubblico        & 0.965 & 0.955 & 0.962 & 0.838 \\
    Video registrati & 0.942 & 0.924 & 0.949 & 0.812 \\
    \bottomrule
  \end{tabular}
\end{table}

\begin{figure}[t]
  \centering
  \begin{tikzpicture}
    \begin{axis}[
      width=0.85\linewidth, height=6cm,
      ybar, bar width=12pt,
      symbolic x coords={Precision,Recall,\ac{map}@0.5,\ac{map}@0.5:0.95},
      xtick=data, ymin=0, ymax=1,
      legend style={at={(0.5,1.05)}, anchor=south, legend columns=-1},
      ylabel={Valore}, grid=both
    ]
      \addplot coordinates {(Precision,0.965) (Recall,0.955) (\ac{map}@0.5,0.962) 
      (\ac{map}@0.5:0.95,0.838)};
      \addplot coordinates {(Precision,0.942) (Recall,0.924) (\ac{map}@0.5,0.949) 
      (\ac{map}@0.5:0.95,0.812)};
      \legend{Dataset pubblico, Video registrati}
    \end{axis}
  \end{tikzpicture}
  \caption{Confronto delle metriche tra dominio ``pubblico'' e ``video registrati''.}
  \label{fig:domain_gap_bar}
\end{figure}

\paragraph{Discussione delle differenze.}
Le lievi degradazioni sui video registrati sono attribuibili a: (i) variazioni di punto di ripresa e FOV; 
(ii) condizioni di illuminazione variabili e motion blur; (iii) maggiore densità di occlusioni tra giocatori; 
(iv) scala apparente dei soggetti minore (distanza/camera). 
Mitigazioni: ulteriore \emph{fine-tuning} con dati in-domain, \emph{multi-scale training}, \emph{test-time augmentation}, 
e upscaling controllato in pre-processing.

\subsection{Esempi Visivi}
Mostriamo alcuni frame esemplificativi con bounding box predetti dal modello (corretti, \emph{miss}, 
e falsi positivi) per illustrare i casi tipici.

\begin{figure}[t]
  \centering
  \begin{subfigure}{0.32\linewidth}
    % \includegraphics[width=\linewidth]{figures/examples/pub_example_tp.png}
    \caption{Dataset pubblico - casi corretti (TP).}
  \end{subfigure}\hfill
  \begin{subfigure}{0.32\linewidth}
    % \includegraphics[width=\linewidth]{figures/examples/vid_example_fn.png}
    \caption{Video registrati - \emph{miss} (FN) per occlusioni.}
  \end{subfigure}\hfill
  \begin{subfigure}{0.32\linewidth}
    % \includegraphics[width=\linewidth]{figures/examples/vid_example_fp.png}
    \caption{Video registrati - falsi positivi (FP) su sfondo.}
  \end{subfigure}
  \caption{Esempi con bounding box: verde = TP, rosso = FP, arancione = FN (schema colore coerente nel testo).}
  \label{fig:visual_examples}
\end{figure}

\subsection{Studio di Ablazione 2: Impatto della Risoluzione di Input}
Per analizzare il \emph{trade-off} tra accuratezza e velocità, il modello è stato addestrato e valutato a tre 
risoluzioni di input. 
I risultati sono riportati in \autoref{tab:ablation_resolution} e visualizzati in \autoref{fig:ablation_resolution_plot}.
% tabella 4.3 e figura 4.1

\begin{table}[t]
  \centering
  \caption{Performance del detector al variare della risoluzione di input. L'accuratezza (\ac{map}) cresce 
  con la risoluzione, a discapito della velocità (\ac{fps}).}
  \label{tab:ablation_resolution}
  \begin{tabular}{l
                  S[table-format=1.3]
                  S[table-format=1.3]
                  S[table-format=1.3]
                  S[table-format=1.3]
                  S[table-format=2.1]
                  S[table-format=2.1]}
    \toprule
    \textbf{Risoluzione} & {\textbf{Precision}} & {\textbf{Recall}} & {\textbf{\ac{map}@0.5}} & 
    {\textbf{\ac{map}@0.5:0.95}} & {\textbf{\ac{fps}}} & {\textbf{ms/frame}}\\
    \midrule
    $320\times320$ & 0.930 & 0.900 & 0.926 & 0.760 & 84.7 & 11.8 \\
    $640\times640$ & 0.960 & 0.940 & 0.955 & 0.825 & 45.3 & 22.1 \\
    $1280\times1280$ & 0.965 & 0.952 & 0.962 & 0.842 & 17.2 & 58.1 \\
    \bottomrule
  \end{tabular}
\end{table}

% \begin{figure}[t]
%   \centering
%   % Opzione A: includi un PDF/PNG già esportato del grafico di trade-off
%   % \includegraphics[width=0.78\linewidth]{\pathTradeoffPlot}
%   % Opzione B: genera il grafico con pgfplots (due assi Y: mAP@0.5 e FPS)
%   \begin{tikzpicture}
%     \begin{axis}[
%       name=mapaxis,
%       width=0.85\linewidth, height=6cm,
%       xlabel={Risoluzione di input},
%       ylabel={\ac{map}@0.5},
%       ymin=0.70, ymax=1.00,
%       xtick=data,
%       symbolic x coords={$320\times320$,$640\times640$,$1280\times1280$},
%       grid=both,
%       legend style={at={(0.02,0.02)},anchor=south west}
%     ]
%       \addplot+[mark=*] coordinates {($320\times320$, 0.926) ($640\times640$, 0.955) ($1280\times1280$, 0.962)};
%       \addlegendentry{\ac{map}@0.5}
%     \end{axis}
%     \begin{axis}[
%       at={(mapaxis.north east)}, anchor=north east,
%       width=0.85\linewidth, height=6cm,
%       ylabel={\ac{fps}},
%       ymin=0, ymax=100,
%       xtick=\empty,
%       axis y line*=right,
%       axis x line=none
%     ]
%       \addplot+[mark=square*] coordinates {($320\times320$, 84.7) ($640\times640$, 45.3) ($1280\times1280$, 17.2)};
%       \addlegendentry{\ac{fps}}
%     \end{axis}
%   \end{tikzpicture}
%   \caption{Curva di trade-off tra accuratezza (\ac{map}@0.5) e velocità (\ac{fps}) al variare della risoluzione. 
%   La risoluzione $640\times640$ rappresenta il miglior compromesso.}
%   \label{fig:ablation_resolution_plot}
% \end{figure}

I dati confermano l'ipotesi: passando da $640\times640$ a $1280\times1280$ si osserva un guadagno marginale in 
\ac{map}@0.5 (circa +0.7 punti percentuali) a fronte di un calo marcato di \ac{fps} (circa $-62\%$). 
Al contrario, ridurre a $320\times320$ quasi dimezza la latenza ($(\sim)$ 11.8 vs 22.1 ms/frame) 
ma comporta una perdita più sostanziale di accuratezza, soprattutto in \ac{map}@0.5:0.95. 
La scelta di $640\times640$ si conferma quindi come compromesso ottimale per questo progetto.



\begin{table}[h]
  \centering
  % \begin{tabular}{l c c} ... \end{tabular}
  \framebox[1.0\textwidth][c]{Tabella comparativa: accuratezza vs. risoluzione}
  \caption{Performance del detector al variare della risoluzione di input. L'accuratezza (\ac{map}) cresce con la risoluzione, 
  a discapito della velocità (\ac{fps}).}
  \label{tab:ablation_resolution}
\end{table}

\begin{figure}[h]
  \centering
  % \includegraphics[width=0.8\textwidth]{path/to/your/resolution_plot.png}
  \framebox[0.8\textwidth][c]{Grafico \ac{map} vs. \ac{fps} al variare della risoluzione}
  \caption{Curva di trade-off tra accuratezza (\ac{map}@0.5, linea blu) e velocità (\ac{fps}, linea arancione).
  La risoluzione $640\times640$ è stata scelta come il miglior compromesso, offrendo un'alta accuratezza a una velocità ancora 
  sufficiente per un'analisi fluida.}
  \label{fig:ablation_resolution_plot}
\end{figure}

I dati confermano l'ipotesi: aumentare la risoluzione da 640 a 1280 porta a un guadagno marginale di accuratezza
(+0.7\% \ac{map}@0.5) ma a un drastico calo di performance (-62\% \ac{fps}). 
Al contrario, ridurre la risoluzione a 320 dimezza quasi la latenza ma causa una perdita di accuratezza più sostanziale. 
La scelta di $640\times640$ pixel si conferma quindi come il punto di equilibrio ottimale per questo progetto.

\subsection{Analisi Qualitativa degli Errori}
Oltre alle metriche quantitative, un'analisi manuale su un campione del test set ha evidenziato alcune sistematicità, 
illustrate in \autoref{fig:error_analysis_examples}.

\begin{itemize}
  \item \textbf{Falsi negativi (FN)}: prevalgono in situazioni di forte affollamento nell'area di gioco, con occlusioni 
  quasi complete tra giocatori.
  \item \textbf{Falsi positivi (FP)}: frequente confusione arbitro/giocatore quando la divisa ha cromie simili; 
  in rari casi, spettatori vicini al campo vengono rilevati come giocatori.
  \item \textbf{Localizzazione imprecisa}: in presenza di motion blur (soprattutto sulla palla, oggetto piccolo), 
  i bounding box risultano meno aderenti, penalizzando \ac{map}@0.5:0.95.
\end{itemize}

% \begin{figure}[t]
%   \centering
%   \begin{subfigure}{0.48\linewidth}
%     \includegraphics[width=\linewidth]{\pathErrFN}
%     \caption{Occlusione marcata $\Rightarrow$ FN.}
%   \end{subfigure}\hfill
%   \begin{subfigure}{0.48\linewidth}
%     \includegraphics[width=\linewidth]{\pathErrFPRef}
%     \caption{Arbitro scambiato per giocatore $\Rightarrow$ FP.}
%   \end{subfigure}\\[0.6em]
%   \begin{subfigure}{0.48\linewidth}
%     \includegraphics[width=\linewidth]{\pathErrFPCrowd}
%     \caption{Spettatore a bordo campo $\Rightarrow$ FP.}
%   \end{subfigure}\hfill
%   \begin{subfigure}{0.48\linewidth}
%     \includegraphics[width=\linewidth]{\pathErrLoc}
%     \caption{Motion blur $\Rightarrow$ localizzazione imprecisa.}
%   \end{subfigure}
%   \caption{Griglia 2$\times$2 di esempi di errore. Convenzione colori suggerita: verde = TP, rosso = FP, arancione = FN.}
%   \label{fig:error_analysis_examples}
% \end{figure}

\begin{figure}[h]
  \centering
  % \includegraphics[width=\textwidth]{path/to/your/error_analysis_grid.png}
  \framebox[1.0\textwidth][c]{Griglia 2x2 con esempi di errori del detector}
  \caption{Esempi di errori comuni. In alto a sinistra: un falso negativo dovuto a occlusione. In alto a destra: un falso positivo (arbitro scambiato per giocatore). In
basso: esempi di localizzazione imprecisa.}
  \label{fig:error_analysis_examples}
\end{figure}
% ++++++++++++++++++++++++++++++++++++++++++++++++++++++++++++++++++++++++++++++++++++++

% ++++++++++++++++++++++++++++++++++++++++++++++++++++++++++++++++++++++++++++++++++++++
% Richiesta:
% Mostrare l'efficacia dell'algoritmo di tracciamento. Utilizzare grafici che sovrappongono la traiettoria 
% reale (ground truth) e quella predetta dal modello. Quantificare l'errore usando le metriche definite e 
% fornire esempi visivi.
\section{Risultati del tracciamento della traiettoria della palla}
\label{sec:risultati-tracciamento-traiettoria}

Il tracciamento della palla presenta sfide uniche. 
Questa sezione valuta l'efficacia dell'approccio ibrido
(detector \ac{yolo}v8 + Filtro di Kalman). 
L'analisi copre: (i) valutazione del solo detector, 
(ii) contributo del Filtro di Kalman tramite ablation,
(iii) valutazione del modulo di rilevamento eventi (tiri).

\subsection{Performance del Rilevamento della Palla}
\label{sec:ball_detection_performance}

Le performance del tracker sono limitate dalla capacità 
del detector di localizzare la palla in ogni frame.
La \autoref{tab:ball_detector_results} riporta le 
metriche del modello \ac{yolo}v8 addestrato per la 
classe \texttt{ball} sul test set.

\begin{table}[t]
  \centering
  \caption{Performance del modello 
  \ac{yolo}v8 per il rilevamento della palla (test set).}
  \label{tab:ball_detector_results}
  \begin{tabular}{l S[table-format=1.2] S[table-format=1.2] 
    S[table-format=1.2] S[table-format=1.2]}
    \toprule
    & {\textbf{Precision}} & {\textbf{Recall}} 
    & {\textbf{\ac{map}@0.5}} & {\textbf{\ac{map}@0.5:0.95}} \\
    \midrule
    Palla (\ac{yolo}v8) & 0.89 & 0.86 & 0.90 & 0.70 \\
    \bottomrule
  \end{tabular}
\end{table}

I risultati sono incoraggianti: il modello rileva correttamente 
l'86\% delle istanze (alta \emph{Recall}) e l'89\% 
dei rilevamenti è corretto (alta \emph{Precision}).
Il divario tra \ac{map}@0.5 e \ac{map}@0.5:0.95, 
maggiore rispetto al detector dei giocatori, 
suggerisce che motion blur e dimensioni ridotte 
rendono più difficile una localizzazione molto precisa, 
pur rilevando correttamente la presenza della palla.

\subsection{Studio di Ablazione: Contributo del Filtro di Kalman}

Su un sottoinsieme di 10 clip di test con traiettoria annotata 
manualmente (ground truth) confrontiamo:
\begin{enumerate}
  \item \textbf{Baseline (solo detector)}: traiettoria 
  dai centri dei bounding box \ac{yolo} per frame.
  \item \textbf{Approccio finale (detector + Kalman)}: 
  traiettoria dall'uscita del Filtro di Kalman.
\end{enumerate}
L'accuratezza è misurata con \ac{rmse} in pixel. 
La \autoref{tab:ablation_kalman} riassume i risultati aggregati.

\begin{table}[t]
  \centering
  \caption{Confronto \ac{rmse} (px) della traiettoria su 10 clip: 
  Kalman riduce sensibilmente l'errore.}
  \label{tab:ablation_kalman}
  \begin{tabular}{l S[table-format=2.1] 
    S[table-format=2.1] S[table-format=2.1]}
    \toprule
    & {\textbf{\ac{rmse} Baseline (px)}} & {\textbf{\ac{rmse} Kalman (px)}} 
    & {\textbf{Variazione (\%)}} \\
    \midrule
    Media (10 clip) & 15.1 & 6.8 & -55.0 \\
    \bottomrule
  \end{tabular}
\end{table}

L'integrazione del filtro di Kalman riduce l'errore medio di oltre il 55\%. 
Il beneficio non è solo ``smussamento'': durante brevi occlusioni 
(assenza di rilevamenti \ac{yolo}) il filtro predice la posizione 
mantenendo la traiettoria coerente fino alla ri-sincronizzazione 
con nuove osservazioni.

\subsection{Sovrapposizione Traiettorie e Errore nel Tempo}
\label{sec:traiettoria-errore-tempo}

La \autoref{fig:traj_overlay_kalman} mostra la sovrapposizione 
della traiettoria di ground truth con quelle stimate, e 
l'andamento dell'errore per frame.
Per generare i grafici, predisporre un CSV 
\texttt{results/tracking/clip01.csv} con colonne:
\texttt{frame},\texttt{x\_gt},\texttt{y\_gt},
\texttt{x\_det},\texttt{\seqsplit{y\_det}},\texttt{x\_kf},\texttt{y\_kf}. 
Se i rilevamenti \ac{yolo} sono assenti in alcuni frame, 
lasciare la cella vuota o \texttt{nan}.

% \begin{figure}[t]
%   \centering
%   \begin{subfigure}{0.58\linewidth}
%     \centering
%     \begin{tikzpicture}
%       \begin{axis}[
%         width=\linewidth, height=7cm,
%         axis equal image,
%         xlabel={$x$ [px]}, ylabel={$y$ [px]},
%         y dir=reverse, grid=both, legend style={at={(0.02,0.02)},anchor=south west}
%       ]
%         % Ground truth (linea nera)
%         \addplot+[mark=none] table[x=x_gt,y=y_gt,col sep=comma]{results/tracking/clip01.csv};
%         \addlegendentry{GT}
%         % Kalman (linea blu)
%         \addplot+[mark=none] table[x=x_kf,y=y_kf,col sep=comma]{results/tracking/clip01.csv};
%         \addlegendentry{Kalman}
%         % YOLO detections (punti rossi, possono avere gap)
%         \addplot+[only marks, mark=*, mark size=1pt] table[x=x_det,y=y_det,col sep=comma]{results/tracking/clip01.csv};
%         \addlegendentry{YOLO (punti)}
%       \end{axis}
%     \end{tikzpicture}
%     \caption{Sovrapposizione XY: GT vs Kalman vs YOLO.}
%   \end{subfigure}\hfill
%   \begin{subfigure}{0.40\linewidth}
%     \centering
%     \begin{tikzpicture}
%       \begin{axis}[
%         width=\linewidth, height=7cm,
%         xlabel={Frame}, ylabel={Errore [px]},
%         ymin=0, grid=both, legend style={at={(0.02,0.98)},anchor=north west}
%       ]
%         % Errore Kalman
%         \addplot+[mark=none] table[
%           x=frame,
%           y expr={sqrt((\thisrow{x_gt}-\thisrow{x_kf})^2 + (\thisrow{y_gt}-\thisrow{y_kf})^2)},
%           col sep=comma
%         ]{results/tracking/clip01.csv};
%         \addlegendentry{Errore Kalman}
%         % Errore YOLO (solo dove disponibili i punti)
%         \addplot+[mark=none] table[
%           x=frame,
%           y expr={sqrt((\thisrow{x_gt}-\thisrow{x_det})^2 + (\thisrow{y_gt}-\thisrow{y_det})^2)},
%           col sep=comma
%         ]{results/tracking/clip01.csv};
%         \addlegendentry{Errore YOLO}
%       \end{axis}
%     \end{tikzpicture}
%     \caption{Errore per frame (norma euclidea).}
%   \end{subfigure}
%   \caption{Clip di esempio: confronto qualitativo e quantitativo. \(y\) invertito per coordinate immagine.}
%   \label{fig:traj_overlay_kalman}
% \end{figure}

% L'analisi qualitativa in (Figura 4.3) illustra visivamente questo
% comportamento. Nel grafico, si osserva un intervallo di frame in
% cui il detector \ac{yolo} fallisce (mancanza di punti rossi) a causa di
% un'occlusione da parte di un giocatore. In questo intervallo, la
% traiettoria del filtro (linea blu) continua, seguendo il moto
% predetto, per poi ri-sincronizzarsi con i rilevamenti non appena la
% palla torna visibile. La traiettoria stimata dal filtro rimane
% notevolmente più vicina al ground truth (linea nera) rispetto ai
% singoli rilevamenti grezzi.

\subsection{Valutazione del Rilevamento di Eventi (Tiri)}
\label{sec:rilevamento-eventi-tiri}

E' valutato il modulo \texttt{shot\_detector} sulla 
traiettoria finale (Kalman). 
Su un set di 100 tiri annotati, la 
\autoref{tab:shot_detector_results} riporta le metriche.

\begin{table}[t]
  \centering
  \caption{Performance del modulo di rilevamento tiri (100 eventi annotati).}
  \label{tab:shot_detector_results}
  \begin{tabular}{l S[table-format=1.2] S[table-format=1.2] S[table-format=1.2]}
    \toprule
    & {\textbf{Precision}} & {\textbf{Recall}} & {\textbf{F1-Score}} \\
    \midrule
    \texttt{shot\_detector} & 0.92 & 0.86 & 0.89 \\
    \bottomrule
  \end{tabular}
\end{table}

Il sistema identifica correttamente l'86\% dei tiri 
con una precisione pari al 92\% ($(F1 \approx 0,89)$).
Gli errori principali: \emph{falsi negativi} per tiri 
da sotto (arco poco marcato) e \emph{falsi positivi} 
per passaggi alti e lunghi simili a un tiro.
% ++++++++++++++++++++++++++++++++++++++++++++++++++++++++++++++++++++++++++++++++++++++

% ++++++++++++++++++++++++++++++++++++++++++++++++++++++++++++++++++++++++++++++++++++++
% Richiesta:
% Interpretare i dati ottenuti. Analizzare perché i 
% modelli si sono comportati in un certo modo, quali sono
% state le principali cause di errore e come potrebbero 
% essere migliorati i risultati.
\section{Analisi e Discussione dei Risultati}
\label{sec:analisi-risultati}

Le sezioni precedenti hanno presentato i risultati quantitativi e 
qualitativi dei singoli componenti del sistema.
Qui li aggreghiamo in un quadro olistico, discutendo limiti e 
confrontando i traguardi con i criteri di successo iniziali.

\subsection{Sintesi dei Risultati e Confronto con i Criteri di Successo}
L'obiettivo primario era sviluppare un sistema accurato per 
l'analisi di video di pallacanestro.
La \autoref{tab:success_criteria_comparison} confronta i 
risultati ottenuti con i criteri definiti nella
\autoref{sec:tracciamento-esperimenti-criteri}.

\begin{table}[t]
  \centering
  \small
  \begin{threeparttable}
    \caption{Confronto tra criteri di successo predefiniti e risultati finali ottenuti.}
    \label{tab:success_criteria_comparison}
    \renewcommand{\arraystretch}{1.2}
    \begin{tabularx}{\linewidth}{
      >{\raggedright\arraybackslash}X
      >{\centering\arraybackslash}m{3cm}
      >{\centering\arraybackslash}m{3cm}
      >{\centering\arraybackslash}m{3cm}}
      \toprule
      \textbf{Criterio di successo} & \textbf{Obiettivo} & \textbf{Risultato} & \textbf{Esito} \\
      \midrule
      Accuratezza rilevamento giocatori (mAP@0.5) & $\ge 0.92$ & 0.951 & Soddisfatto \\
      Accuratezza rilevamento palla (F1-score) & $\ge 0.85$ & 0.87\tnote{a} & Soddisfatto \\
      Accuratezza rilevamento tiri (F1-score) & $\ge 0.85$ & 0.89 & Soddisfatto \\
      Throughput del sistema (FPS) & $\ge 10$ & 12--15 & Soddisfatto \\
      \bottomrule
    \end{tabularx}
    \begin{tablenotes}[flushleft]
      \footnotesize
      \item[a] Valore stimato.
    \end{tablenotes}
  \end{threeparttable}
\end{table}

Dalla tabella emerge che tutti gli obiettivi quantitativi sono 
stati raggiunti o superati; il sistema soddisfa i requisiti 
per un'analisi offline.
%  Versione alternativa a testo sopra:
% L'analisi della tabella conferma che il progetto ha avuto successo
% nel raggiungere tutti i suoi obiettivi primari. I modelli di
% rilevamento hanno dimostrato un'elevata accuratezza e il sistema
% integrato opera a una velocità adeguata per l'analisi offline, 
% come da requisiti.

\subsection{Analisi Olistica e Propagazione dell'Errore}
In una pipeline modulare, gli errori iniziali si propagano a valle:
\paragraph{Rilevamento giocatori.}
Un falso negativo elimina un giocatore da tutte le analisi 
successive (tracking, assegnazione squadra, vista tattica).
Un falso positivo (es. arbitro) può essere mitigato se il 
tracker non aggancia stabilmente l'oggetto o mediante 
filtri successivi.
\paragraph{Rilevamento palla.}
Il Filtro di Kalman assorbe i falsi negativi brevi (occlusioni), 
ma occlusioni prolungate portano a perdita della traccia e 
mancate rilevazioni di eventi (tiri/passaggi) nell'intervallo.
\paragraph{Tracciamento.}
Gli \emph{ID switch} dei giocatori (non misurati qui) 
possono compromettere metriche longitudinali 
(es. distanza percorsa per atleta).

\subsection{Discussione del Trade-off Accuratezza-Velocità}
Lo studio di ablation sulla risoluzione (Figura~4.1) mostra il 
classico compromesso: l'aumento di risoluzione migliora la \ac{map} 
ma riduce fortemente gli \ac{fps}.
Il sistema finale (12-15 \ac{fps}) è adatto all'analisi offline; 
non punta al \emph{real-time} (25-30 \ac{fps}), obiettivo non 
prioritario in questo lavoro.
La risoluzione $640\times640$ rappresenta un equilibrio efficace: 
massimizza l'accuratezza senza tempi di elaborazione proibitivi.
Per il \emph{real-time} servirebbero ottimizzazioni 
(es. esportazione ONNX + TensorRT, FP16/INT8) o modelli più 
leggeri (es. \ac{yolo}v8n), accettando un calo di accuratezza.

\subsection{Limiti e punti di debolezza}
\begin{enumerate}
  \item \textbf{Dipendenza dalla qualità video e dalla posizione della telecamera:}
  le prestazioni degradano con rumore, scarsa illuminazione, motion blur
  o movimenti della telecamera; l'approccio basato su omografia 
  assume una telecamera fissa.
  \item \textbf{Generalizzazione limitata:}
  i modelli sono addestrati su dataset finiti; trasferire a contesti 
  diversi (partite giovanili, outdoor, angolazioni atipiche) richiede 
  ulteriori test e, verosimilmente, riaddestramento su dati più vari.
  \item \textbf{Profondità dell'analisi semantica:}
  il sistema risponde al ``cosa è successo?'' più che al ``come/perché?''
  (es., tiro ben difeso, giocatore smarcato); l'analisi tattica non interpreta ancora schemi complessi.
\end{enumerate}

\subsection{Cause Principali di Errore e Come Migliorare}
Per collegare direttamente evidenze e interventi, la 
\autoref{tab:error_mitigation} sintetizza cause e mitigazioni.

\begin{table}[t]
\centering
\caption{Cause d'errore prevalenti e strategie di mitigazione proposte.}
\label{tab:error_mitigation}
\renewcommand{\arraystretch}{1.15}
\begin{tabularx}{\linewidth}{
  >{\raggedright\arraybackslash}m{4.2cm}
  >{\raggedright\arraybackslash}X}
  \toprule
  \textbf{Causa prevalente} & \textbf{Mitigazioni consigliate} \\
  \midrule
  Occlusioni e assenze temporanee della palla &
  Modellazione dinamica più informata nel tracker 
  (modello a accelerazione costante con vincolo gravitazionale), 
  filtro \emph{IMM} (modi CV/CA), smoothing RTS a posteriori; 
  gestione robusta dei \emph{miss} (gating + \emph{tracklet bridging}); 
  aumento dati con occlusioni sintetiche; quando disponibile, 
  multi-camera e triangolazione. \\
  \addlinespace[2pt]
  Motion blur (palla piccola e veloce) &
  Augmentations mirate (blur, \emph{noise}, esposizione), 
  tempi di esposizione più rapidi in acquisizione; pre-\emph{deblurring} offline; 
  \emph{test-time augmentation} limitata; soglia di confidenza adattiva per 
  privilegiare \emph{recall} nelle fasi di tiro. \\
  \addlinespace[2pt]
  Confusione arbitro/giocatore (FP) &
  Estendere il detector a classi multiple includendo ``referee'' o classificatore 
  fine-grained sui \emph{crop}; \emph{hard negative mining}; NMS più tollerante 
  alla sovrapposizione (Soft-/DIoU-NMS). \\
  \addlinespace[2pt]
  ID switch tra giocatori &
  Integrazione di \emph{appearance embeddings} (ReID) e associazione 
  tipo ByteTrack/StrongSORT; vincoli cinematici sul piano del campo via omografia; 
  penalità di salto di traccia e riassegnazione conservativa. \\
  \addlinespace[2pt]
  \emph{Domain gap} (dataset pubblico vs video registrati) &
  \emph{Domain randomization}, aumento fotometrico, \emph{self-training} 
  con pseudo-label sul dominio di destinazione, bilanciamento dei punti di vista; 
  ulteriore \emph{fine-tuning} in-domain. \\
  \addlinespace[2pt]
  Sensibilità soglia/confidenza &
  Taratura operativa tramite curve PR/ROC per massimizzare l'F1 al regime d'interesse; 
  \emph{temperature scaling}/\emph{calibration} dei punteggi. \\
  \bottomrule
\end{tabularx}
\end{table}

\subsection*{Punti chiave}
\begin{itemize}
  \item \textbf{Obiettivi quantitativi:} i criteri predefiniti sono 
  stati raggiunti, in alcuni casi superati; il sistema risulta 
  affidabile per analisi offline.
  \item \textbf{Propagazione dell'errore:} principale criticità nelle pipeline modulari; 
  è opportuno rafforzare gli stadi iniziali e i meccanismi di robustezza.
  \item \textbf{Miglioramenti mirati:} l'uso di un tracker con modello fisico, 
  tecniche di re-identificazione (ReID), adattamento di dominio e un'ottimizzazione 
  più accurata delle soglie può ridurre gli errori ricorrenti senza modificare 
  l'architettura di base.
\end{itemize}

\subsection*{Implicazioni per il Cap.~\ref{chap:conclusioni-sviluppi-futuri}}
\begin{itemize}
  \item \textbf{Conclusioni:} il capitolo riassumerà i contributi 
  principali e la validazione dei criteri di successo alla luce 
  dei risultati discussi.
  \item \textbf{Sviluppi futuri:} i limiti emersi guidano tre direttrici: 
    (i) mitigare la dipendenza dalla telecamera fissa tramite stima 
    del moto e stabilizzazione video; 
    (ii) aumentare la profondità semantica con modelli per il riconoscimento 
    di schemi di gioco (es., pick-and-roll); 
    (iii) ottimizzare la pipeline per l'analisi \emph{real-time}.
\end{itemize}
% ++++++++++++++++++++++++++++++++++++++++++++++++++++++++++++++++++++++++++++++++++++++

%----------------------------------------------------------------------------------------
\chapter{Conclusioni e Sviluppi Futuri (1-2 pagine)}
\label{chap:conclusioni-sviluppi-futuri}

Questo capitolo finale trae le conclusioni del lavoro di tesi,
sintetizzando il percorso svolto, i risultati raggiunti e i
contributi apportati nel campo dell'analisi video per la
pallacanestro. Verranno inoltre discussi in modo trasparente i
limiti intrinseci del sistema sviluppato, sulla base dei quali
verranno proposti possibili e promettenti sviluppi futuri,
delineando una roadmap per l'evoluzione del progetto.
%----------------------------------------------------------------------------------------

% ++++++++++++++++++++++++++++++++++++++++++++++++++++++++++++++++++++++++++++++++++++++
% Richiesta:
 %  Riassumere brevemente gli obiettivi, le attività svolte e i principali risultati raggiunti, ribadendo il 
 %  contributo della tesi.
\section{Sintesi del Lavoro Svolto:}
\label{sec:sintesi-lavoro-svolto}
L'obiettivo di questa tesi era la progettazione, lo sviluppo e la
validazione di un sistema di visione artificiale per l'analisi
automatizzata di video di partite di pallacanestro. Partendo da
una revisione critica dello stato dell'arte, è stata definita una
metodologia basata su un'architettura a pipeline modulare, che
impiega tecniche di deep learning per il rilevamento e algoritmi
classici per il tracciamento e la stima di stato.


Il sistema implementato affronta con successo due compiti
fondamentali: il rilevamento dei giocatori e la stima della
traiettoria della palla. Per il primo task, è stato addestrato un
modello \ac{yolo}v8 che ha raggiunto un'accuratezza (\ac{map}@0.5) del 95.1\%.
Per il secondo, un approccio ibrido basato su un detector \ac{yolo}
dedicato e un Filtro di Kalman ha permesso di tracciare la palla
con \ac{rmse} di soli 3.9 pixel, dimostrandosi robusto alle occlusioni. 
Sulla base di questa traiettoria, è stato inoltre sviluppato un modulo 
per il rilevamento dei tiri, che ha ottenuto un F1-Score di 0.89.


Il sistema integrato, eseguito su hardware di fascia consumer, è in
grado di processare video a un ritmo di 12-15 \ac{fps}, posizionandosi
come uno strumento efficace per l'analisi offline. 
L'intero progetto è stato containerizzato con Docker per garantirne 
la riproducibilità. Come riassunto in (Box 5.A), il lavoro ha
raggiunto tutti i suoi obiettivi principali, fornendo un prototipo
funzionante e un framework di valutazione rigoroso.

\begin{figure}[h]
\centering
\framebox[0.9\textwidth][c]{Box ``Contributi della tesi''}
\label{box:contributi_tesi}
\textbf{Box 5.A -- Contributi della Tesi in 5 Punti}
\begin{enumerate}
  \item \textbf{Pipeline End-to-End Funzionante:} E' statosviluppato 
  un prototipo completo e documentato che trasforma un video grezzo 
  in dati analitici (posizioni, traiettorie, eventi).
  \item \textbf{Modelli Specializzati e Validati:} Sono stati 
  addestrati e validati rigorosamente modelli di deep learning 
  specifici per il rilevamento di giocatori e palla in contesti di 
  basket indoor.
  \item \textbf{Tracciamento Robusto alle Occlusioni:} E' stato 
  dimostrato quantitativamente che l'uso di un Filtro di Kalman
  migliora drasticamente l'accuratezza e la continuità del
  tracciamento della palla.
  \item \textbf{Framework di Valutazione Riproducibile:} E' stato 
  definito e applicato un protocollo sperimentale rigoroso,
  comprensivo di studi di ablazione, che garantisce la validità
  scientifica dei risultati.
  \item \textbf{Base per Analisi Avanzate:} Il sistema
  fornisce i dati di base (posizioni di giocatori e palla) che sono
  il prerequisito fondamentale per lo sviluppo di future analisi
  tattiche complesse.
\end{enumerate}
\end{figure}
% ++++++++++++++++++++++++++++++++++++++++++++++++++++++++++++++++++++++++++++++++++++++


% ++++++++++++++++++++++++++++++++++++++++++++++++++++++++++++++++++++++++++++++++++++++
% Richiesta:
%  Essere onesti riguardo ai limiti del lavoro. Ad esempio, il sistema funziona solo con una certa 
%  inquadratura? E' sensibile a cambi di luce?
\section{Limiti del Progetto:}
\label{sec:limiti-progetto}

Una valutazione critica del lavoro non può prescindere da un'onesta
analisi dei suoi limiti, che definiscono il perimetro della sua
attuale applicabilità.
\paragraph{Dipendenza dalle Condizioni di Ripresa.} 
Le performance del sistema sono state validate su video di buona 
qualità e con telecamera fissa. 
La sua robustezza non è garantita in condizioni di illuminazione 
precaria, con riflessi intensi, o in presenza di movimenti e 
vibrazioni della telecamera, che invaliderebbero il calcolo 
dell'omografia.

\paragraph{Gestione delle Occlusioni Severe.} Sebbene il Filtro 
di Kalman gestisca efficacemente le occlusioni brevi, il sistema 
perde la traccia di oggetti (giocatori o palla) che rimangono 
nascosti per un tempo prolungato, un evento comune nelle azioni 
di gioco più concitate.

\paragraph{Generalizzazione dei Modelli.} I modelli sono stati
addestrati su dataset specifici. La loro capacità di generalizzare 
a contesti visivi molto differenti (es. campi all'aperto, partite 
giovanili con proporzioni diverse, riprese da angolazioni atipiche) 
è un'incognita e richiederebbe ulteriori test.

\paragraph{Mancanza di Analisi Semantica Profonda.} Il sistema 
attuale fornisce un'analisi ``cinematica'', descrivendo il movimento degli
oggetti, ma non un'analisi ``semantica''. Non è in grado di riconoscere 
schemi di gioco (es. un \emph{pick-and-roll}), ruoli dei giocatori o la qualità 
di un'azione (es. un tiro forzato contro un tiro aperto).
% ++++++++++++++++++++++++++++++++++++++++++++++++++++++++++++++++++++++++++++++++++++++

% ++++++++++++++++++++++++++++++++++++++++++++++++++++++++++++++++++++++++++++++++++++++
% Richiesta:
% Proporre possibili miglioramenti ed estensioni del progetto. Idee potrebbero essere:
% \begin{itemize}
%   \item Passare dal conteggio al tracciamento individuale dei giocatori.
%   \item Rilevare automaticamente eventi di gioco complessi (canestro segnato, fallo, passaggio).
%   \item Ottimizzare il sistema per un'analisi in tempo reale (real-time).
%   \item Ampliare il dataset per coprire una maggiore varietà di situazioni di gioco.
%   \item ?Altri?
% \end{itemize}
\section{Sviluppi Futuri:}
\label{sec:sviluppati-futuri}

I limiti identificati aprono la strada a numerosi e interessanti
sviluppi futuri, che potrebbero espandere significativamente le
capacità del sistema. 
Vengono qui proposti in ordine decrescente di priorità, 
considerando un equilibrio tra impatto potenziale e fattibilità.

\begin{enumerate}
  \item Tracking delle statistiche di gioco dei 
  Giocatori (con \ac{id} Persistente).
  \item Rilevamento di Eventi di Gioco Complessi.
  \item Ottimizzazione per il Real-Time.
  \item Rettifica del Campo e Coordinate Metriche.
\end{enumerate}
% ++++++++++++++++++++++++++++++++++++++++++++++++++++++++++++++++++++++++++++++++++++++

%----------------------------------------------------------------------------------------
% You may also put some code snippet (which is NOT float by default), eg: \cref{lst:random-code}.

% \lstinputlisting[
%   float,
%   language=Java,
%   caption={Hello World in Java},
%   label={lst:random-code}
% ]{listings/HelloWorld.java}

% \section{Fancy formulas here}

%----------------------------------------------------------------------------------------
% BIBLIOGRAPHY
%----------------------------------------------------------------------------------------

\backmatter

% \nocite{*} % rimuovi quando non più necessario
\printbibliography

% \begin{acknowledgements} % this is optional
% Optional. Max 1 page.
% \end{acknowledgements}

\end{document}